\subsection{\texorpdfstring{The relation \(\sum_{i} e_i f_i \leqslant n\)}{The relation Σ ei fi ≤ n}}

Let K be a field, v a valuation on K and L a finite extension of K of degree n.~Let \((v_{1}^{\prime}, \ldots, v_{r}^{\prime})\) be valuations on L extending v no two of which are equivalent; if they are independent (which is always the case if v is of height 1), then

\[
\sum_ {i = 1} ^ {r} e (v _ {i} ^ {\prime} / v) f (v _ {i} ^ {\prime} / v) \leqslant n
\]

by Proposition 2 (formulae (6) and (7)). We shall see that this result is true in the general case. To be precise:

THEOREM 1. Let K be a field, v a valuation on K and L a finite extension of K of degree n.~Then:

\begin{enumerate}
\def\labelenumi{(\alph{enumi})}
\item
  Every complete system \((v_i')_{i\in I}\) of extensions of \(v\) to \(\mathbf{L}\) is finite.
\item
  \(\sum_{i\in I}e(v_i^{\prime}|v)f(v_i^{\prime}|v)\leqslant n\) and a fortiori Card(I) \(\leqslant n.\)
\item
  No two of the rings of the \(v_{i}'\) are comparable with respect to inclusion.
\end{enumerate}

Since the theorem is trivial if \(v\) is improper, we shall assume that \(v\) is not improper. Let \((v_1', \ldots, v_s')\) be any finite family of valuations on \(L\) extending \(v\) , no two of which are equivalent. We shall first prove that \(\sum_{i=1}^{s} e(v_i'/v)f(v_i'/v) \leqslant n\) . This will prove (a) and (b).

We argue by induction on s and suppose therefore that the inequality has been established for the case of 0, 1, \ldots, s - 1 valuations. We distinguish two cases.

\begin{enumerate}
\def\labelenumi{(\arabic{enumi})}
\tightlist
\item
  Suppose that there exist at least two independent valuations \(v_i'\) . Then there exists (§ 7, no. 2, Remark 1) a partition \([1, s] = I_1 \cup \cdots \cup I_t\) of \([1, s]\) such that:
\end{enumerate}

\begin{enumerate}
\def\labelenumi{(\roman{enumi})}
\item
  for \(v_{i}^{\prime}\) and \(v_{j}^{\prime}\) to be dependent, it is necessary and sufficient that \(i\) and \(j\) belong to the same \(I_{k}\) ;
\item
  \(\operatorname{Card}(\mathbf{I}_k) < s\) for all \(k\) .
\end{enumerate}

We choose in each \(I_{k}\) an index \(i(k)\) . Let \(\hat{L}_{i(k)}\) denote the completion of L with respect to \(v_{i(k)}^{\prime}\) and \(n(k) = [\hat{L}_{i(k)} : \hat{K}]\) . For all \(i \in I_{k}\) , \(v_{i}^{\prime}\) defines on L the same topology as \(v_{i(k)}^{\prime}\) (§ 7, no. 2, Proposition 3) and hence may be extended to a valuation \(\hat{v}_{i}^{\prime}\) on \(\hat{L}_{i(k)}\) whose restriction to \(\hat{K}\) is \(\hat{v}\) . Since no two of the \(v_{i}^{\prime}\) for \(i \in I_{k}\) are equivalent, the same is true of the \(\hat{v}_{i}^{\prime}\) . The induction hypothesis applied to the ordered pair ( \(\hat{K}, \hat{L}_{i(k)}\) ) shows, by virtue of Proposition 2 (a), formula (4), that \(\sum_{i\in I_k}e(v_i^{\prime}|v)f(v_i^{\prime}|v)\leqslant n(k)\) . As \(\sum_{k = 1}^{t}n(k)\leqslant n\) (Proposition 2 (b), formula

(7)), certainly \(\sum_{i=1}^{s} e(v_i'/v) f(v_i'/v) \leqslant n\) .

\begin{enumerate}
\def\labelenumi{(\arabic{enumi})}
\setcounter{enumi}{1}
\tightlist
\item
  We now pass to the case where any two of the \(v_i'\) are dependent. Let \(A_i'\) be the ring of \(v_i' (1 \leqslant i \leqslant s)\) ; writing A for the ring of \(v, A_i' \cap K = A\) for all \(i\) . Let \(B'\) be the subring of L generated by \(A_1', \ldots, A_s'\) ; we write \(B = B' \cap K\) ; then \(B \supset A\) . Then B is the ring of a valuation \(w\) on K and \(B'\) the ring of a valuation \(w'\) which is not improper and extends \(w\) ( \(\S 7\) , no. 2, Proposition 4); the field \(\kappa(B')\) is an extension of \(\kappa(B)\) of degree \(f(w'/w)\) . Consider the canonical images \(\bar{A}_i', \bar{A}\) of \(A_i'\) and A in \(\kappa(B')\) ; then \(\bar{A}\) is the ring of a valuation \(\bar{v}\) on \(\kappa(B)\) and the \(\bar{A}_i'\) are rings of valuations \(\bar{v}_i'\) on \(\kappa(B')\) extending \(\bar{v}\) . As the \(A_i'\) generate \(B'\) , the \(\bar{A}_i'\) generate \(\kappa(B')\) and hence the \(\bar{v}_i'\) are not all dependent ( \(\S 7\) , no. 2, Proposition 4). From the first part of the proof,
\end{enumerate}

\[
\sum_ {i = 1} ^ {s} e \left(\bar {v} _ {i} ^ {\prime} / \bar {v}\right) f \left(\bar {v} _ {i} ^ {\prime} / \bar {v}\right) \leqslant \left[ \kappa \left(\mathbf {B} ^ {\prime}\right): \kappa (\mathbf {B}) \right] = f \left(w ^ {\prime} / w\right)
\]

and hence

\[
\sum_ {i = 1} ^ {s} e (w ^ {\prime} / w) e (\bar {v} _ {i} ^ {\prime} / \bar {v}) f (\bar {v} _ {i} ^ {\prime} / \bar {v}) \leqslant e (w ^ {\prime} / w) f (w ^ {\prime} / w) \leqslant n \quad \text {(no. 1, Lemma 1)}.
\]

The proof of (a) and (b) will therefore be completed if we prove that

\[
f (\bar {v} _ {i} ^ {\prime} / \bar {v}) = f (\bar {v} _ {i} ^ {\prime} / \bar {v}), \quad e (w ^ {\prime} / w) e (\bar {v} _ {i} ^ {\prime} / \bar {v}) = e (v _ {i} ^ {\prime} / v).\tag{8}
\]

For this, we note that v and \(\bar{v}\) (resp. \(v_{i}^{\prime}\) and \(\bar{v}_{i}^{\prime}\) ) have the same residue field (§ 4, no. 1, Corollary to Proposition 2); this proves the first equation. For the second there is, by virtue of the Remark in § 4, no. 3, the following commutative diagram, where the rows are exact sequences and the vertical arrows represent canonical injections:

\[
\begin{array}{c} 0 \to \Gamma_ {\bar {v}} \to \Gamma_ {v} \to \Gamma_ {w} \to 0 \\ \downarrow \quad \downarrow \quad \downarrow \\ 0 \to \Gamma_ {\bar {v} _ {i} ^ {\prime}} \to \Gamma_ {v _ {i} ^ {\prime}} \to \Gamma_ {w ^ {\prime}} \to 0 \end{array}
\]

We deduce that there is an exact sequence

\[
0 \rightarrow \Gamma_ {\bar {v} _ {i} ^ {\prime}} / \Gamma_ {\bar {v}} \rightarrow \Gamma_ {v _ {i} ^ {\prime}} / \Gamma_ {v} \rightarrow \Gamma_ {w ^ {\prime}} / \Gamma_ {w} \rightarrow 0
\]

by Chapter I, § 1, no. 4, Proposition 2, which proves the second formula (8). To complete the proof of Theorem 1, it remains to prove (c). If the ring of \(v_i'\) contains that of \(v_j'\) , \(\Gamma_{v_i'}\) is identified with a quotient group \(\Gamma_{v_j'} / H\) , \(H\) being an isolated subgroup (§ 4, no. 3). As the composite canonical mapping

\[
\Gamma_ {v} \rightarrow \Gamma_ {v _ {j} ^ {\prime}} \rightarrow \Gamma_ {v _ {j} ^ {\prime}} / H = \Gamma_ {v _ {i} ^ {\prime}}
\]

is injective, \(H \cap \Gamma_{v} = \{0\}\) , whence \(H = \{0\}\) (Lemma 3, no. 1). Then \(v_{i}'\) and \(v_{j}'\) are equivalent, whence i = j.

Remark. The intersection C of the rings \(A_{i}^{\prime}\) of the valuations \(v_{i}^{\prime}\) ( \(i \in I\) ) is the integral closure of A in L ( \(\S 1\) , no. 3, Corollary 3 to Theorem 3); it follows moreover from (c) and \(\S 7\) , no. 1, Propositions 1 and 2 that C is a semi-local ring, that its maximal ideals are the intersections \(m_{i} = C \cap m(A_{i}^{\prime})\) and that \(A_{i}^{\prime} = C_{m}\) for all \(i \in I\) .

\subsection{Initial ramification index}

DEFINITION 4. Let G be a totally ordered commutative group and H a subgroup of G of finite index. The number of major subsets of G consisting of strictly positive elements and containing all the elements \textgreater0 of H is called the initial index of H in G and denoted by \(\varepsilon(\mathbf{G}, \mathbf{H})\) .

This initial index is a natural number by virtue of the following proposition:

PROPOSITION 3. Under the hypotheses of Definition 4, if the set of strictly positive elements of G has no least element, then \(\varepsilon(G, H) = 1\) for all H. If there exists a least element \(>0\) of G and G' denotes the subgroup it generates, then

\[
\varepsilon (\mathrm{G}, \mathrm{H}) = (\mathrm{G} ^ {\prime}: (\mathrm{G} ^ {\prime} \cap \mathrm{H})).
\]

In the first case, let x be an element \textgreater0 in G. The set of \(y \in G\) such that 0 \textless{} y \textless{} x is infinite and hence there exist two elements of this set which are distinct and congruent modulo H; their difference is an element z of H such that 0 \textless{} z \textless{} x. Hence every major subset which contains all the strictly positive elements of H contains x and hence all the elements \textgreater0 of G.

In the second case, let \(x\) be the least element \(>0\) of \(G\) and let \(n\) be the least integer \(>0\) such that \(nx \in H\) . Clearly \(n = (G': (G' \cap H))\) . On the other hand, writing \(M(y)\) for the set of \(z \in G\) such that \(y \leqslant z\) ( \(y \in G\) ), it is immediately seen that the major sets of Definition 4 are just \(M(x), M(2x), \ldots, M(nx)\) .

COROLLARY. The initial index \(\varepsilon(\mathbf{G},\mathbf{H})\) divides the index (G: H) and is equal to it if G is isomorphic to Z.

In particular, \(\varepsilon(G, H) \leqslant (G: H)\) .

DEFINITION 5. Let K be a field, L a finite extension of K, w a valuation on L, v its restriction to K and \(\Gamma_{w}\) and \(\Gamma_{v}\) their order groups. The initial index of \(\Gamma_{v}\) in \(\Gamma_{w}\) is called the initial ramification index of w with respect to v (or with respect to K) and denoted by \(\varepsilon(w/v)\) .

From the above corollary, \(\varepsilon(w/v)\) divides \(e(w/v)\) with equality in the case of a discrete valuation.

PROPOSITION 4. Under the hypotheses of Definition 5, let A and m (resp. A' and m') be the ring and ideal of the valuation v (resp. w). Then

\[
[ \mathrm{A} ^ {\prime} / \mathrm{mA} ^ {\prime}: \mathrm{A} / \mathrm{m} ] = \varepsilon (w / v) f (w / v).
\]

The ideals of A' containing mA' and distinct from A' correspond to the major subsets of \(\Gamma_{w}\) consisting of elements \textgreater0 and containing the elements \textgreater0 of \(\Gamma_{v}\) ( \(\S3\) , no. 5, Corollary to Proposition 7). They are therefore equal in number to \(\varepsilon(w/v)\) and, as they form a totally ordered set under inclusion, this number is equal to the length of the quotient ring A'/mA'. Now a module of length 1 over A' is a 1-dimensional vector space over A'/m' and hence a module of length \(f(w/v)\) over A; hence, as A'/mA' is of length \(\varepsilon(w/v)\) over A', it is of length \(\varepsilon(w/v)f(w/v)\) over A, that is over A/m.

\subsection{\texorpdfstring{The relation \(\sum_{i} e_{i} f_{i} = n\)}{The relation Σ ei fi = n}}

PROPOSITION 5. Let K be a field, v a valuation on K, A its ring, m its ideal, L a finite extension of K of degree n, B the integral closure of A in L and \((v_{i}^{\prime})_{1 \leq i \leq s}\) a complete system of extensions of v to L. Then

\[
[ \mathrm{B} / \mathrm{mB}: \mathrm{A} / \mathrm{m} ] = \sum_ {i = 1} ^ {s} \varepsilon (v _ {i} ^ {\prime} / v) f (v _ {i} ^ {\prime} / v).
\]

Let \(A_{i}\) be the ring of \(v_{i}^{\prime}\) ; then \(A_{i} = B_{m_{i}}\) , where \(m_{i}\) runs through the family of maximal ideals of B (no. 3, Remark). Let \(q_{i}\) be the saturation of mB with respect to \(m_{i}\) (Chapter II, § 2, no. 4). By Chapter V, Corollary 3 to Proposition 1, no. 1, § 2, the canonical homomorphism \(B/mB \rightarrow \prod_{i=1}^{s} B/q_i\) is an isomorphism and \(m_i\) is the only maximal ideal of B containing \(q_i\) . Hence \(B/q_i\) is canonically isomorphic to \((\mathrm{B}/q_i)_{m_i}\) (Chapter II, §3, no. 3, Proposition 8), that is to

\[
\mathrm{B} _ {\mathrm{m} _ {i}} / \mathrm{mB} _ {\mathrm{m} _ {i}} = \mathrm{A} _ {i} / \mathrm{mA} _ {i}.
\]

Therefore there is a canonical isomorphism \(B/mB \rightarrow \prod_{i=1}^{s} A_{i}/mA_{i}\) , whence the result by virtue of Proposition 4 of no. 4.

COROLLARY. With the same hypotheses and notation,

\[
[ \mathrm{B} / \mathrm{mB}: \mathrm{A} / \mathrm{m} ] = \sum_ {i = 1} ^ {s} \varepsilon (v _ {i} ^ {\prime} / v) f (v _ {i} ^ {\prime} / v) \leqslant \sum_ {i = 1} ^ {s} e (v _ {i} ^ {\prime} / v) f (v _ {i} ^ {\prime} / v) \leqslant n.
\]

We know that \(\varepsilon(v_i'/v) \leqslant e(v_i'/v)\) (no. 4, Corollary to Proposition 3) and \(\sum_{i=1}^{s} e(v_i'/v)f(v_i'/v) \leqslant n\) (no. 3, Theorem 1).

THEOREM 2. With the hypotheses and notation of Proposition 5, the following conditions are equivalent:

\begin{enumerate}
\def\labelenumi{(\alph{enumi})}
\item
  B is a finitely generated A-module;
\item
  B is a free A-module;
\item
  \([\mathbf{B} / \mathfrak{m}\mathbf{B}\colon \mathbf{A} / \mathfrak{m}] = n;\)
\end{enumerate}

\[
\text {(d)} \sum_ {i = 1} ^ {n} e (v _ {i} ^ {\prime} / v) f (v _ {i} ^ {\prime} / v) = n \text { and } \varepsilon (v _ {i} ^ {\prime} / v) = e (v _ {i} ^ {\prime} / v) \text { for   all } i.
\]

The equivalence of (a) and (b) follows from Lemma 1, § 3, no. 6. Clearly (b) implies (c) (Algebra, Chapter II, § 1, no. 5, formula (19)). The equivalence of (c) and (d) follows from the Corollary to Proposition 5. It remains to show that (c) implies (b).

In general, if M is an A-module, we shall denote by V(M) the vector space M/mM over A/m. Hypothesis (c) means that \(\dim(\mathrm{V}(\mathrm{B})) = n\) . Let \(x_{1}, \ldots, x_{n}\) be elements of B whose canonical images in V(B) form a basis of V(B) and let \(L \subset B\) be the sub-A-module which they generate. As L is torsion-free and finitely generated, it is free (§ 3, no. 6, Lemma 1). We shall see that B = L. Let \(y \in B\) ; we write \(M = L + Ay\) ; this is also a free A-module. The canonical injections \(L \to M \to B\) give canonical homomorphisms \(\mathrm{V}(\mathrm{L}) \to \mathrm{V}(\mathrm{M}) \to \mathrm{V}(\mathrm{B})\) . As the ranks of L and M are \(\leqslant n\) , so are the dimensions of V(L) and V(M). Now, by hypothesis, \(\mathrm{V}(\mathrm{L}) \to \mathrm{V}(\mathrm{B})\) is surjective and V(B) is n-dimensional hence V(L) and V(M) are n-dimensional and \(\mathrm{V}(\mathrm{L}) \to \mathrm{V}(\mathrm{M})\) is surjective. As M is finitely generated, \(L \to M\) is surjective (Chapter II, § 3, no. 2, Corollary 1 to Proposition 4), whence \(L = M, y \in L\) and B = L. Hence B is free.

Remark (1) If \(v\) is discrete, \(\varepsilon(v_i'/v) = e(v_i'/v)\) (no. 4) and condition (d) reduces to \(\sum_{i=1}^{s} e(v_i'/v)f(v_i'/v) = n\) .

COROLLARY 1. With the same hypotheses and notation, suppose further that v is discrete and L separable. Then

\[
\sum_ {i = 1} ^ {s} e (v _ {i} ^ {\prime} / v) f (v _ {i} ^ {\prime} / v) = n.
\]

The integral closure B of A is then a free A-module of rank n, since A is a principal ideal domain (Chapter V, § 1, no. 6, Corollary 2 to Proposition 18).

COROLLARY 2. Let K be a field, v a discrete valuation on K with respect to which K is complete and L a finite extension of K of degree n.~Then v admits a unique (up to equivalence) extension v' to L, the ring A' of v' is a finitely generated free module over the ring A of v and \(e(v'/v)f(v'/v)=n\) .

All the extensions of \(v\) to \(L\) are dependent (no. 2, Corollary to Proposition 2); since they are discrete (no. 1, Corollary 3 to Proposition 1), they are therefore equivalent (§ 4, no. 5, Proposition 6 (c)). This shows the uniqueness of \(v'\) . The integral closure of \(A\) in \(L\) is therefore \(A'\) (§ 1, no. 3, Corollary 3 to Theorem 3). As \(v\) is discrete, the topology induced on \(A\) by that on \(K\) is the m-adic topology (where \(m = m(A)\) ); the ring \(A\) is complete, for it is closed in \(K\) . We conclude that, since \(A'/mA'\) is a finite-dimensional vector \((A/m)\) -space (no. 4, Proposition 4), \(A'\) is a finitely generated A-module (Chapter III, § 2, no. 9, Corollary 3 to Proposition 12). It is therefore free and \(e(v'/v)f(v'/v) = n\) by Theorem 2.

COROLLARY 3. Suppose that \(v\) is of height 1 and that the equivalent conditions of Theorem 2 hold; if \(\hat{\mathbf{L}}_i\) is the completion of \(\mathbf{L}\) with respect to \(v_i'\) , the degree \(n_i = [\hat{\mathbf{L}}_i: \hat{\mathbf{K}}]\) is equal to \(e(v_i'|v)f(v_i'|v)\) for all \(i\) and the canonical homomorphism

\[
\phi \colon \hat {\mathbf {K}} \otimes_ {\mathbf {K}} \mathbf {L} \rightarrow \prod_ {i = 1} ^ {s} \hat {\mathbf {L}} _ {i}
\]

(no. 2, Proposition 2). is bijective. For all \(x \in \mathbf{L}\) , the characteristic polynomial \(\mathrm{Pc}_{\mathrm{L} / \mathbb{K}}(x; \mathbf{X})\) is equal to the product of the characteristic polynomials \(\mathrm{Pc}_{\hat{\mathrm{L}}_i / \hat{\mathbf{K}}}(x; \mathbf{X})\) ( \(1 \leqslant i \leqslant s\) ); in particular,

\[
\left\{ \begin{array} { l } \mathrm{Tr} _ { \mathbf { L } / \mathbb { K } } ( x ) = \sum _ { i = 1 } ^ { s } \mathrm{Tr} _ { \hat { \mathbf { L } } _ { i } / \hat { \mathbf { K } } } ( x ) \\ \mathrm{N} _ { \mathbf { L } / \mathbb { K } } ( x ) = \prod _ { i = 1 } ^ { s } \mathrm{N} _ { \hat { \mathbf { L } } _ { i } / \hat { \mathbf { K } } } ( x ) \\ v ( \mathrm{N} _ { \mathbf { L } / \mathbb { K } } ( x ) ) = \sum _ { i = 1 } ^ { s } n _ { i } v _ { i } ^ { \prime } ( x ) . \\ \end{array} \right.\tag{9}
\]

(The last relation in (9) is meaningful, for we may obviously assume that the \(v_{i}^{\prime}\) , which are of height 1 by Corollary 2 to Proposition 1 of no. 1, take, as does \(v\) , their values in a subgroup of \(\mathbf{R}\) .)

As no two of the \(v_{i}^{\prime}\) are equivalent and they are of height 1, they are independent and Proposition 2 of no. 2 shows therefore that \(e(v_{i}^{\prime}/v)f(v_{i}^{\prime}/v) \leqslant n_{i}\) for all \(i\) and \(\sum_{i=1}^{s} n_i \leqslant n\) . The first assertion therefore follows from these inequalities and the relation \(\sum_{i=1}^{s} e(v_i'/v) f(v_i'/v) = n\) . Under the isomorphism \(\phi\) the endomorphism \(z \mapsto z(1 \otimes x)\) of \(\hat{\mathbf{K}} \otimes_{\mathbb{K}} \mathbf{L}\) (for \(x \in \mathbf{L}\) ) is transformed into the endomorphism of \(\prod_{i=1}^{s}\hat{L}_i\) leaving invariant each of the factors and reducing on each factor to multiplication by \(x\) (L being canonically imbedded in its completion \(\hat{L}_i\) ); whence the assertion relating to the characteristic polynomial of \(x\) and the first two formulae of (9). Finally, let E be a finite quasi-Galois extension of \(\hat{K}\) , containing \(\hat{L}_i\) ; as \(\hat{K}\) is complete and \(\hat{v}\) of height 1, there exists only one valuation (up to equivalence) \(w\) on E extending \(\hat{v}\) (no. 2, Corollary 1 to Proposition 2); then, for every \(\hat{K}\) -automorphism \(\sigma\) of E, \(w(\sigma(x)) = v_i'(x)\) . Therefore

\[
\hat {v} \left(\mathrm{N} _ {\hat {\mathrm{L}} _ {i} / \hat {\mathrm{K}}} (x)\right) = n _ {i} v _ {i} ^ {\prime} (x)
\]

(Algebra, Chapter VIII, § 12, no. 2, formula (15)), which proves the formula of (9).

COROLLARY 4. Under the hypotheses of Corollary 3, if L is a separable extension of K, each of the \(\hat{\mathbf{L}}_i\) is a separable extension of \(\hat{\mathbf{K}}\) . If further L is a Galois extension of K with Galois group \(\mathcal{G}\) and \(\mathcal{G}_i\) denotes the decomposition group of the ideal of \(v_i'\) in B (Chapter V, § 2, no. 2, Definition 2), then \(\hat{\mathbf{L}}_i\) is a Galois extension of \(\hat{\mathbf{K}}\) whose Galois group is isomorphic to \(\mathcal{G}_i\) .

Clearly \(\hat{\mathbf{L}}_i = \hat{\mathbf{K}}(\mathbf{L})\) ; hence, if \(\mathbf{L}\) is separable over \(\mathbf{K}\) , \(\hat{\mathbf{L}}_i\) is separable over \(\hat{\mathbf{K}}\) (Algebra, Chapter V, § 7, no. 6, Proposition 10). Suppose now that \(\mathbf{L}\) is Galois. Every automorphism \(\sigma \in \mathcal{G}_i\) is continuous on \(\mathbf{L}\) with the topology defined by \(v_i'\) , the fact that no two of the ideals of the \(v_i'\) are comparable with respect to inclusion (§ 7, no. 2, Corollary 1 to Theorem 1) necessarily implying that \(v_i' = v_i' \circ \sigma\) by definition of \(\mathcal{G}_i\) ; hence \(\sigma\) may be extended by continuity to a \(\hat{\mathbf{K}}\) -automorphism \(\hat{\sigma}\) of \(\hat{\mathbf{L}}_i\) . This proves that the number of \(\hat{\mathbf{K}}\) -automorphisms of \(\hat{\mathbf{L}}_i\) is at least equal to \(\operatorname{Card}(\mathcal{G}_i)\) . But as the valuations \(v_i'\) are pairwise conjugate under \(\mathcal{G}\) (Chapter V, § 2, no. 3, Proposition 6), \(s = (\mathcal{G}: \mathcal{G}_i)\) , whence

\[
\operatorname{Card} \left(\mathcal {G} _ {i}\right) = n / s \leqslant n.
\]

and on the other hand \(n = sn_{i}\) by Corollary 3; this proves that \(\hat{\mathbf{L}}_{i}\) is a Galois extension of \(\hat{\mathbf{K}}\) and that the extensions by continuity of the automorphisms \(\sigma \in \mathcal{G}_{i}\) are the only \(\hat{\mathbf{K}}\) -automorphisms of \(\hat{\mathbf{L}}_{i}\) .

Remark (2). Part of the above results extends to the case of valuations on a not necessarily commutative field K (cf.~§ 3, no. 1). Let L be an extension field of K and let \(v'\) be a valuation on L, \(v\) its restriction to K and A' and A the respective rings of the valuations \(v'\) and v; then there is defined a ramification index \(e(v'/v)\) as in no. 1; on the other hand, \(\kappa(\mathbf{A})\) is identified with a subfield of \(\kappa(\mathbf{A}')\) and the (left) residue rank of \(v'\) with respect to v is defined to be the number \(f(v'/v)\) equal to the dimension of the left vector \(\kappa(\mathbf{A})\) -space \(\kappa(\mathbf{A}')\) , if this dimension is finite, and \(+\infty\) in the opposite case. Then, if L is a left vector K-space of finite dimension n, Lemma 2 of no. 1 and its proof go over unchanged. Moreover, if K is complete with respect to v, the assertions of Corollary 2 to Theorem 2 of no. 5 (other than the existence of \(v'\) ) are also valid (n denoting the dimension of L as a left vector K-space) with the following proof:

In the first place the topology defined by \(v'\) on \(L\) is Hausdorff and compatible with its left vector \(K\) -space structure and hence two extensions of \(v\) to \(L\) give the same topology on \(L\) (§ 5, no. 2, Proposition 4), which proves that these extensions are the same up to equivalence (§ 6, no. 2). We show next that, if \(m = m(A)\) , \(A'/mA'\) is a left vector \((A/m)\) -space of dimension \(e(v'/v)f(v'/v)\) . Writing \(e = e(v'/v)\) , we may assume that \(v(K^*) = Z\) and \(v'(L^*) = e^{-1}Z\) ; let \(u'\) be an element of \(L\) such that \(v'(u') = e^{-1}\) and \(u\) an element of \(K\) such that \(v(u) = 1\) ; hence \(u = zu'^e\) , where \(z \in L\) is such that \(v'(z) = 0\) . As \(m\) is generated by \(u\) (as a left or right ideal of \(A\) ), \(mA' = u'^eA' = A'u'^e\) and it suffices to prove that, for \(0 \leqslant k \leqslant e - 1\) , \(A'u'^k/A'u'^{k+1}\) is a left vector \((A/m)\) -space of dimension \(f(v'/v)\) . But \(t \mapsto tu'^k\) is an isomorphism of the left A-module \(A'\) onto the left A-module \(A'u'^k\) mapping \(A'u'\) to \(A'u'^{k+1}\) and which therefore gives by taking quotients an \((A/m)\) -isomorphism of \(A'/A'u'\) onto \(A'u'^k/A'u'^{k+1}\) , whence our assertion by definition of \(f(v'/v)\) , \(u'\) generating the maximal ideal of \(A'\) . The proof is completed as when \(K\) and \(L\) are commutative (the fact that a finitely generated torsion-free A-module is free being proved as in § 3, no. 6, Lemma 1).

\subsection{Valuation rings in an algebraic extension}

PROPOSITION 6. Let K be a field, v a valuation on K, A its ring, L an algebraic extension of K and A' the integral closure of A in L. Let V be the set of valuation rings on L which extend v and M' the set of maximal ideals of A'. Then the mapping V \(\mapsto\) m(V) \(\cap\) A' is a bijection of V onto M' and \(m' \mapsto A'_{m'}\) is the inverse bijection.

Every maximal ideal \(m'\) of \(A'\) is such that \(m' \cap A\) is the maximal ideal of A (Chapter V, §2, no. 1, Proposition 1) and \(A_{m'}'\) is dominated by a valuation ring V of L (which is therefore the ring of a valuation on L extending v) (§1, no. 2, Corollary to Theorem 2). The field L is the union of a directed family of sub-extensions \(K_{\alpha}\) of L which are finite over K and it will suffice, in order to see that \(V = A_{m'}'\) , to prove that \(V \cap K_{\alpha} = A_{m'}' \cap K_{\alpha}\) for all \(\alpha\) . Now, if we write \(A_{\alpha}' = A' \cap K_{\alpha}\) , \(A_{\alpha}'\) is the integral closure of A in \(K_{\alpha}\) and hence is the intersection of the rings of the valuations on \(K_{\alpha}\) which extend v and these rings \(V_{i\alpha}\) are finite in number and are the local rings \((\mathrm{A}_{\alpha}')_{\mathrm{m}'_{i\alpha}}\) of \(A_{\alpha}' (1 \leqslant i \leqslant n)\) , where the \(m'_{i\alpha}\) are the distinct maximal ideals of \(A_{\alpha}' (no. 3, Remark)\) ; but \(m' \cap A_{\alpha}'\) is one of the \(m'_{i\alpha}\) and \(V \cap K_{\alpha}\) is therefore equal to the corresponding local ring \((\mathrm{A}'_{\alpha})_{m'_{i\alpha}} \subset \mathrm{A}'_{m'}\) , which completes the proof that \(V = A'_{m'}\) . Conversely, if \(V \in V\) , then \(A' \subset V\) (§3, no. 3, Proposition 6) and, if \(m' = m(V) \cap A'\) , then \(m' \cap A = m\) , hence \(m'\) is a maximal ideal of \(A'\) (Chapter V, §2, no. 1, Proposition 1) and the above argument shows that \(V = A'_{m'}\) .

PROPOSITION 7. Let K be a field, L a quasi-Galois extension of K and f and \(f'\) places of L with values in the same field F. Suppose that the restrictions of f and \(f'\) to K coincide. Then there exists a K-automorphism s of L such that \(f' = f \circ s\) .

Let A be the ring of the place of K the common restriction of f and \(f'\) . The rings of f and \(f'\) contain the integral closure \(A'\) of A in L (§ 1, no. 3, Corollary 3 to Theorem 3) and hence (Chapter V, § 2, no. 3, Corollary 1 to Proposition 6) there exists a K-automorphism s of L such that the restrictions of \(f'\) and \(f \circ s\) to \(A'\) are equal; if \(m'\) is the common kernel of these restrictions, \(m' \cap A\) is the maximal ideal of A, hence \(m'\) is a maximal ideal of \(A'\) and the places \(f'\) and \(f \circ s\) coincide on the ring \(A_{m'}\) ; but by Proposition 6 the only valuation ring of L dominating \(A_{m'}\) is the ring \(A_{m'}\) itself and hence the rings of the places \(f'\) and \(f \circ s\) are the same.

COROLLARY 1. Let K be a field, v a valuation on K, L a quasi-Galois extension of K and v' and v'' two extensions of v to L. Then there exists a K-automorphism s of L such that v'' is equivalent to v' o s.

Let \(f'\) and \(f''\) be the places of K associated with \(v'\) and \(v''\) ; replacing them if need be by equivalent places, it may be assumed that they both take their values in the algebraic closure of the residue field of \(v\) (no. 1, Proposition 1). Then there exists a K-automorphism \(s\) of L such that \(f'' = f' \circ s\) (Proposition 7); thus \(v''\) is equivalent to \(v' \circ s\) by virtue of the correspondence between places and valuations (§ 3, no. 3).

COROLLARY 2. Let K be a field, f a place of K (resp. v a valuation on K) and L a radicial extension of K. Then all the extensions of f (resp. v) to L are equivalent.

L is a quasi-Galois extension and its only automorphism is the identity. Corollary 2 therefore follows from Proposition 7 (resp. Corollary 1).

PROPOSITION 8. Let K be a field, v a valuation on K, L a finite quasi-Galois extension of K of degree n and \((v_{i}^{\prime})_{1\leq i\leq g}\) a complete system of extensions of v to L. Then \(e(v_{i}^{\prime}|v)\) and \(f(v_{i}^{\prime}|v)\) have values e and f independent of i. Then \(efg \leqslant n\) . If the integral closure in L of the ring A of v is a finitely generated A-module, then \(efg = n\) .

This follows immediately from Theorems 1 (no. 3) and 2 (no. 5).

\subsection{The extension of absolute values}

PROPOSITION 9. Let K be a field, L an algebraic extension of K and f an absolute value on K. Then f can be extended to an absolute value on L.

Suppose first that there exists a valuation \(v\) on \(\mathbf{K}\) with real values such that \(f(x) = e^{-v(x)}\) . There exists a valuation \(v'\) on \(\mathbf{L}\) whose restriction to \(\mathbf{K}\) is equivalent to \(v\) (§ 3, no. 3, Proposition 5). Then \(v'\) is of height 0 or 1 (no. 1, Corollary 2 to Proposition 1) and therefore may be assumed to have real values. The restriction of the mapping \(x \mapsto e^{-v'(x)}\) to \(\mathbf{K}\) is an absolute value equivalent to \(f\) and hence of the form \(f^s\) with \(s > 0\) (General Topology, Chapter IX, § 3, no. 2, Proposition 5). We conclude that

\[
x \mapsto e ^ {- v ^ {\prime} (x) / s}
\]

is an absolute value on \(\mathbf{L}\) extending \(f\) .

Suppose now that \(f\) is not ultrametric. Then \(K\) is identified with a subfield of \(\mathbf{C}\) such that \(f(x) = |x|^s\) where \(0 \leqslant s \leqslant 1\) (§ 6, no. 4, Theorem 2). As \(\mathbf{C}\) is algebraically closed, \(L\) is identified with a subfield of \(\mathbf{C}\) and the absolute value \(x \mapsto |x|^s\) extends \(f\) .

PROPOSITION 10. Let K be a field, f an absolute value on K such that K is complete and not discrete with respect to f and L an algebraic extension of K. Then f can be extended uniquely to an absolute value \(f'\) on L and, if L is of finite degree n, then

\[
f ^ {\prime} (x) = \left(f (\mathrm{N} _ {\mathrm{L/K}} (x))\right) ^ {1 / n}
\]

for all \(x \in \mathbf{L}\) .

The existence of \(f'\) follows from Proposition 9 and its uniqueness (over every finite sub-extension of \(L\) and therefore over the whole of \(L\) ) from Lemma 2 of §6, no. 4. Let \(f'\) be the unique extension of \(f\) to the algebraic closure of \(K\) and

suppose L is of finite degree n.~We know that \(\mathrm{N}_{\mathrm{L}/\mathrm{K}}(x)=\prod_{i=1}^{n} x_i\) , where each \(x_{i}\) is a conjugate of x over K (Algebra, Chapter VIII, § 12, no. 2, Proposition 4). In view of the uniqueness of \(f', f'(x_{i})=f'(x)\) for all i, whence the stated formula.

PROPOSITION 11. Let K be a field, f a non-ultrametric absolute value on K, \(\hat{K}\) the completion of K with respect to f, \(\hat{f}\) the continuous extension of f to \(\hat{K}\) and L a finite extension of K of degree n.

\begin{enumerate}
\def\labelenumi{(\alph{enumi})}
\item
  Let \(f'\) be an absolute value on \(\mathbf{L}\) extending \(f\) ; let \(\hat{\mathbf{L}}_{f'}\) denote the completion of \(\mathbf{L}\) with respect to \(f'\) and let \(\hat{\mathbf{K}}\) be identified with the closure of \(\mathbf{K}\) in \(\hat{\mathbf{L}}_{f'}\) ; then \([\hat{\mathbf{L}}_{f'}: \hat{\mathbf{K}}] \leqslant n\) .
\item
  The absolute values on L extending f are finite in number. If they are denoted by \(f_{1}^{\prime}, \ldots, f_{s}^{\prime}\) and the completion of L with respect to \(f_{i}^{\prime}\) by \(\hat{L}_{i}\) , the canonical mapping
\end{enumerate}

\(\hat{\mathbf{K}}\otimes_{\mathbf{K}}\mathbf{L}\rightarrow \prod_{i = 1}^{n}\hat{\mathbf{L}}_i\) is an isomorphism and

\[
\sum_ {i = 1} ^ {s} \left[ \hat {\mathrm{L}} _ {i}: \hat {\mathrm{K}} \right] = n.\tag{10}
\]

The proof is the same as that for the analogous assertions in Proposition 2 (no. 2). The references

§ 7, no. 2, Theorem 1; § 5, no. 2, Corollary to Proposition 4

should be replaced by the following

§ 7, no. 3, Theorem 2; Topological Vector Spaces, Chapter I,

§ 2, no. 3, Corollary 1 to Theorem 2.

Observe that two extensions of \(f\) to \(\mathbf{L}\) which define the same topology are equal (General Topology, Chapter IX, § 3, no. 2, Proposition 5). Finally, as \(f\) is not ultrametric, \(\mathbf{K}\) is of characteristic 0 and hence the Jacobson radical of \(\hat{\mathbf{K}} \otimes_{\mathbf{K}} \mathbf{L}\) is zero.

\textbf{Remarks}\par

\begin{enumerate}
\def\labelenumi{(\arabic{enumi})}
\item
  Proposition 11 (b) shows that every composite extension of \(\hat{\mathbf{K}}\) and \(\mathbf{L}\) over \(\mathbf{K}\) is isomorphic to one of the completions \(\hat{\mathbf{L}}_i\) and that these are composite extensions no two of which are isomorphic.
\item
  We know that the completions \(\hat{\mathbf{K}}\) and \(\hat{\mathbf{L}}_i\) are isomorphic to \(\mathbf{R}\) or \(\mathbf{C}\) (§ 6, no. 4, Theorem 2). If \(\hat{\mathbf{K}}\) is isomorphic to \(\mathbf{C}\) , so is \(\hat{\mathbf{L}}_i\) for all \(i\) and (10) shows that the number of extensions \(f_i'\) is equal to \(n\) . If \(\hat{\mathbf{K}}\) is isomorphic to \(\mathbf{R}\) (for example if \(\mathbf{K} = \mathbf{Q}\) ), let \(r_1\) (resp. \(r_2\) ) denote the number of indices \(i\) such that \(\hat{\mathbf{L}}_i\) is isomorphic to \(\mathbf{R}\) (resp. \(\mathbf{C}\) ); then (10) may be written:
\end{enumerate}

\[
r _ {1} + 2 r _ {2} = n.\tag{11}
\]

PROPOSITION 12. Let K be a field, f an absolute value on K, L a quasi-Galois extension of K and \(f'\) and \(f''\) two extensions of f to L. Then there exists a K-automorphism s of L such that \(f'' = f' \circ s\) .

If \(f\) is ultrametric, Corollary 1 to Proposition 7 (no. 6) shows that there exists a K-automorphism \(s\) of \(L\) such that \(f''\) and \(f' \circ s\) are equivalent absolute values; then there exists a real number \(a > 0\) such that \(f''(x) = (f'(s(x)))^a\) for all \(x \in L\) . If \(f\) is not improper, take \(x \in K^*\) such that \(f(x) \neq 1\) , which shows that \(a = 1\) . If \(f\) is improper, so are \(f'\) and \(f''\) (Corollary 2 to Proposition 1, no. 1) and \(s\) may be taken to be the identity automorphism.

If \(f\) is not ultrametric, there exist \(\mathbf{Q}\) -isomorphisms \(u', u''\) of \(\mathbf{L}\) onto subfields of \(\mathbf{C}\) and real exponents \(a' > 0\) , \(a'' > 0\) such that \(f'(x) = |u'(x)|^{a'}\) and

\[
f ^ {\prime \prime} (x) = | u ^ {\prime \prime} (x) | ^ {a ^ {\prime \prime}}
\]

for all \(x \in \mathbf{L}\) (§ 6, no. 4, Theorem 2). Taking \(x = 2\) , it is seen that \(a' = a''\) . The restrictions of \(u'\) and \(u''\) to \(\mathbf{K}\) extend by continuity to isomorphisms \(u_1\) and \(u_2\) of \(\hat{\mathbf{K}}\) onto \(\mathbf{R}\) (resp. \(\mathbf{C}\) ). Then \(u_2 \circ \bar{u}_1^{-1}\) is an automorphism of the valued field \(\mathbf{R}\) (resp. \(\mathbf{C}\) ) and is therefore the identity (resp. the identity or the automorphism \(c: \zeta \to \bar{\zeta}\) ). Replacing if need be \(u'\) by \(c \circ u'\) , it is seen that the restrictions of \(u'\) and \(u''\) to \(\mathbf{K}\) may be assumed to coincide. Identifying \(\mathbf{K}\) with a subfield of \(\mathbf{C}\) by means of this common restriction, \(u'\) and \(u''\) are K-isomorphisms of \(\mathbf{L}\) onto subfields of \(\mathbf{C}\) . As \(\mathbf{L}\) is a quasi-Galois extension of \(\mathbf{K}\) , there exists a K-automorphism \(s\) of \(\mathbf{L}\) such that \(u'' = u' \circ s\) ; since \(a' = a''\) , we deduce immediately that \(f'' = f' \circ s\) .

Remark (3). If \(\hat{\mathbf{K}}\) is isomorphic to \(\mathbf{R}\) , Proposition 12 shows that all the completions \(\hat{\mathbf{L}}_t\) of \(\mathbf{L}\) (in the notation of Proposition 11) are isomorphic to one another. Thus, with the notation of Remark 2 above, either \(r_1 = n\) and \(r_2 = 0\) , or \(r_1 = 0\) and \(2r_2 = n\) .

\section{APPLICATION: LOCALLY COMPACT FIELDS}

\subsection{The modulus function on a locally compact field}

Let K be a locally compact field (not necessarily commutative). Recall that the function mod (or mod \(_K\) ) has been defined (Integration, Chapter VII, § 1, no. 10, Definition 6) on K as follows: mod \(_K(0) = 0\) and for \(x \neq 0\) in K, the number mod \(_K(x)\) is the modulus of the automorphism \(y \mapsto xy\) of the additive group of K.

PROPOSITION 1. If K is a locally compact field, the function mod \(_{K}\) belongs to \(\mathcal{V}(K)\) ( \(\S\) 6, no. 1). Moreover:

\begin{enumerate}
\def\labelenumi{(\roman{enumi})}
\item
  If \(s > 0\) is such that \((\mathrm{mod}_{\mathbf{K}})^s = g\) is an absolute value, then \(g\) defines the topology on \(\mathbf{K}\) .
\item
  If K is not discrete and \(mod_{K}\) is an ultrametric absolute value, there exists a normed discrete valuation v on K whose ring is compact and whose residue field is finite with q elements, so that \(mod_{K} = q^{-v}\) . The topology on K is defined by v.
\end{enumerate}

This follows from § 6, no. 1, Proposition 1, § 5, no. 1, Proposition 2 and Integration, Chapter VII, § 1, no. 10, Propositions 12 and 13.

PROPOSITION 2. Let K, K' be two (not necessarily commutative) locally compact fields such that K is a topological subfield of K' and K is not discrete. Then:

\begin{enumerate}
\def\labelenumi{(\roman{enumi})}
\item
  \(\mathbf{K}'\) is a finite dimensional left (resp. right) vector space over \(\mathbf{K}\) .
\item
  If \(\mathbf{K}\) is contained in the centre of \(\mathbf{K}'\) , then, for all \(x \in \mathbf{K}'\) .
\end{enumerate}

\[
\operatorname{mod} _ {\mathbf {K} ^ {\prime}} (x) = \operatorname{mod} _ {\mathbf {K}} (\mathrm{N} _ {\mathbf {K} ^ {\prime} / \mathbf {K}} (x)).\tag{1}
\]

As K is a complete valued field which is not discrete, assertion (i) follows from Topological Vector Spaces, Chapter I, §2, no. 4, Theorem 3; assertion (ii) is just Integration, Chapter VII, §1, no. 11, Proposition 17.

COROLLARY 1. Every locally compact field whose centre is not discrete is of finite rank over its centre.

The centre Z of a locally compact field K is closed in K and therefore locally compact.

COROLLARY 2. Let K' be a locally compact field and K a closed subfield of K' (neither necessarily commutative). If K' is a left (resp. right) vector space of finite dimension n over K, then

\[
\operatorname{mod} _ {\mathrm{K} ^ {\prime}} (x) = (\operatorname{mod} _ {\mathrm{K}} (x)) ^ {n} \quad \text {for all} \quad x \in \mathrm{K}.\tag{2}
\]

In general it is known that in a (left or right) vector space of finite dimension n over K, the homothety with ratio \(x \in K\) has modulus equal to \((\text{mod}_{\mathbf{K}}(x))^{n}\) ; it suffices to apply this to \(K'\) .

\subsection{Existence of representatives}

PROPOSITION 3. Let K be a (not necessarily commutative) locally compact field which is not discrete and whose topology is defined by a discrete valuation v; let A be the ring and m the ideal of v and let us write \(\operatorname{Card}(\mathbf{A}/\mathfrak{m}) = q = p^{f}\) (p prime). Then, there exists a system of representatives S of A/m in A and a uniformizer u for v such that \(0 \in S\) , \(S^{*} = S \cap K^{*}\) is a cyclic subgroup of \(K^{*}\) and \(u^{-1} Su = S\) . Moreover, every element of A may be written uniquely in the form \(\sum_{i=0}^{\infty} s_{i} u^{i}\) , where \(s_{i} \in S\) .

We shall use the following lemma:

LEMMA 1. Let \(x, y\) be two permutable elements of \(A\) such that \(x - y \in \mathfrak{m}^j\) ( \(j \geqslant 1\) ); then \(x^{p^n} - y^{p^n} \in \mathfrak{m}^{j + n}\) for every integer \(n \geqslant 0\) .

By induction on \(n\) , it is reduced to proving the lemma for \(n = 1\) . Then

\[
x ^ {p} - y ^ {p} = (x - y) \left(x ^ {p - 1} + x ^ {p - 2} y + \dots + y ^ {p - 1}\right);
\]

the second factor is a sum of \(p\) terms congruent to one another mod. \(m\) and, as \(A / m\) is of characteristic \(p, p\) . \(1 \in m\) in \(A\) and hence

\[
x ^ {p - 1} + x ^ {p - 2} y + \dots + y ^ {p - 1} \in \mathfrak {m};
\]

whence \(x^{p} - y^{p}\in \mathfrak{m}^{j + 1}\)

We know that the multiplicative group \((\mathbf{A} / \mathfrak{m})^*\) is a cyclic group with \(q - 1\) elements (Algebra, Chapter V, § 11, no. 1, Theorem 1); let \(x\) be a representative in A of a generator of this group; then \(x^{q} - x \in \mathfrak{m}\) , whence, by Lemma 1, \(x^{q^{n + 1}} - x^{q^n} \in \mathfrak{m}^{1 + fn}\) , since \(x^{q}\) and \(x\) are permutable. This proves that \((x^{q^n})_{n \geqslant 0}\) is a Cauchy sequence in A; as A is compact and hence complete, this sequence has a limit \(s\) in A which obviously satisfies \(s \equiv x\) (mod. \(\mathfrak{m}\) ) and \(s^q = s\) . As \(s \neq 0\) , \(s^{q - 1} = 1\) , more precisely \(s\) is a primitive \((q - 1)\) -th root of unity in A. Clearly the set S consisting of 0 and the powers \(s^j\) ( \(0 \leqslant j \leqslant q - 2\) ) is a system of representatives of the classes of A mod. \(\mathfrak{m}\) and is invariant under multiplication in A.

Now let \(a\) be a uniformizer for \(v\) and consider the inner automorphism \(y \mapsto a^{-1}ya\) of \(K\) ; it maps \(A\) to itself, \(m\) to itself and therefore, taking quotients, it defines an automorphism of the field \(A/m\) ; it is known (Algebra, Chapter V, § 11, no. 4, Proposition 5) that such an automorphism is of the form \(z \mapsto z^{p^r}\) , where \(0 \leqslant r \leqslant f - 1\) . Then \(a^{-1}s^j a \equiv s^{jp^r} \pmod{m}\) for \(0 \leqslant j \leqslant q - 2\) ; as \(a \in m\) and \(s \notin m\) , this implies that \(s^{-j}as^{jp^r} \equiv a \pmod{m^2}\) .

Let us write

\[
u = \sum_ {j = 0} ^ {q - 2} s ^ {- j} a s ^ {j p ^ {r}}.
\]

Then \(u \equiv (q - 1)a \equiv -a \pmod{m^2}\) since \(p.1 \in m\) ; we conclude that \(u\) is also a uniformizer for \(v\) ; moreover

\[
s ^ {- 1} u s ^ {p ^ {r}} = u\tag{3}
\]

whence we deduce that \(u^{-1} \mathrm{S}u = \mathrm{S}\) .

Finally, for all \(x \in \mathbf{A}\) there exists a unique sequence \((s_i)\) ( \(i \in \mathbf{N}\) ) such that \(s_i \in \mathbf{S}\) for all \(i\) and \(x \equiv \sum_{i=0}^{n} s_i u^i \pmod{\mathfrak{m}^{n+1}}\) for all \(n \geqslant 0\) : it is immediate by induction on \(n\) , every element \(t\) of \(\mathfrak{m}^{n+1}\) satisfying a relation of the form \(t \equiv t'u^{n+1} \pmod{\mathfrak{m}^{n+2}}\) , where \(t'\) is a uniquely determined element of \(S\) . Then \(x = \sum_{i=0}^{\infty} s_i u^i\) and the family \((s_i)\) satisfying this relation and such that \(s_i \in S\) for all \(i\) is determined uniquely.

\subsection{Structure of locally compact fields}

The completions R and \(Q_{p}\) of the field Q with respect to the non-improper absolute values on Q (p any prime) are locally compact. On the other hand, for every power \(q = p^{f}\) of a prime number p, the field \(\mathbf{F}_{q}((\mathrm{T}))\) of formal power series over the finite field \(F_{q}\) with the valuation defined in §3, no. 4, Example 3 is locally compact: for the maximal ideal of the valuation ring \(F_{q}[[T]]\) is generated by T; we know that this ring is complete with the (T)-adic topology (Chapter III, §2, no. 6, Proposition 6) and, as the residue field \(F_{q}\) is finite, Proposition 2 of §5, no. 1 proves our assertion. Conversely:

THEOREM 1. Let K be a (not necessarily commutative) locally compact field which is not discrete.

\begin{enumerate}
\def\labelenumi{(\roman{enumi})}
\item
  If \(\mathbf{K}\) is of characteristic 0 and \(\mathrm{mod}_{\mathbf{K}}\) is not an ultrametric absolute value then \(\mathbf{K}\) is isomorphic to one of the fields \(\mathbf{R}\) , \(\mathbf{C}\) or \(\mathbf{H}\) .
\item
  If \(\mathbf{K}\) is of characteristic 0 and \(\mathrm{mod}_{\mathbf{K}}\) is an ultrametric absolute value, \(\mathbf{K}\) is an algebra of finite rank over a \(p\) -adic field \(\mathbf{Q}_p\) .
\item
  If \(\mathbf{K}\) is of characteristic \(p \neq 0\) , it is isomorphic to a field with centre a field of formal power series \(\mathbf{F}_q((\mathrm{T}))\) (where \(q\) is a power of \(p\) ) and of finite rank over its centre.
\end{enumerate}

\begin{enumerate}
\def\labelenumi{(\roman{enumi})}
\item
  It follows from Ostrowski's Theorem (§ 6, no. 4, Theorem 2) that K is a topological field isomorphic to an everywhere dense subfield of R, C or H and, as K is complete, it is isomorphic to R, C or H.
\item
  Let A be the ring of the absolute value \(mod_{K}\) and m its maximal ideal. We know that A/m is a finite field (§ 5, no. 1, Proposition 2) and hence the absolute value induced by \(mod_{K}\) on Q has a finite residue field, which is only possible if it is equivalent to a \(p\) -adic absolute value (§ 6, no. 3, Proposition 4); the closure of \(\mathbf{Q}\) in \(\mathbf{K}\) is therefore isomorphic to \(\mathbf{Q}_p\) and is contained in the centre of \(\mathbf{K}\) since the latter is closed in \(\mathbf{K}\) ; we conclude using Proposition 2 of no. 1.
\item
  The second assertion follows from the first and the Corollary to Proposition 2 of no. 2. To show the first assertion, note that \(\mathrm{mod}_{\mathbf{K}}\) is necessarily an ultrametric absolute value (§ 6, no. 2, Corollary to Proposition 3); in the notation of the proof of Proposition 3 of no. 2, the centre \(Z\) of \(K\) consists of the elements which commute with both \(s\) and \(u\) ; but by virtue of (3),
\end{enumerate}

\[
u ^ {- q} s u ^ {q} = s ^ {q p ^ {r}} = s,
\]

so that \(u^q \in \mathbf{Z}\) and we conclude that \(\mathbf{Z}\) is not discrete. As \(\mathbf{Z}\) is locally compact, we may confine our attention to the case where \(\mathbf{K}\) is commutative. The sub- \(\mathbf{F}_p\) -algebra \(\mathbf{F}_p[s]\) in \(\mathbf{K}\) is then a finite field since \(s^{q-1} = 1\) and obviously \(y^q = y\) for every element of this field, which is therefore identical with \(S\) and isomorphic to \(\mathbf{F}_q\) since \(S \subset \mathbf{F}_p[s]\) has \(q\) elements. Since the sum of two elements of \(S\) is in \(S\) , the mapping which maps each formal power series \(\sum_{i=0}^{\infty} s_i T^i \in \mathbf{F}_q[[T]]\) to the element \(\sum_{i=0}^{\infty}s_{i}u^{i}\) is a bijective homomorphism of the ring \(F_{q}[[T]]\) onto the ring A, whence immediately the conclusion.

COROLLARY 1. Every locally compact field which is not discrete is of finite rank over its centre.

COROLLARY 2. Every locally compact field is connected or totally disconnected; if it is connected, it is isomorphic to R, C or H.

If the topology on a field K is defined by an ultrametric absolute value, K is totally disconnected with this topology.

Remark. Let \(s\) be an integer \(>0\) ; the subfield \(\mathbf{F}_{q}((\mathrm{T}^{s})) = \mathbf{L}\) of \(\mathbf{K} = \mathbf{F}_{q}((\mathrm{T}))\) is closed in \(\mathbf{K}\) and \(e(\mathrm{K/L}) = s\) and \(f(\mathrm{K/L}) = 1\) . It is therefore seen that there are closed subfields \(\mathbf{L}\) of \(\mathbf{K}\) which are not discrete and such that \(e(\mathrm{K/L})\) (and a fortiori the degree {[}K: L{]}) is arbitrarily large (contrary to what happens for locally compact fields of characteristic 0, where every locally compact subfield \(\mathbf{L}\) of such a field \(\mathbf{K}\) necessarily contains \(\mathbf{R}\) or \(\mathbf{Q}_{p}\) and therefore {[}K: L{]} is bounded).

\section{EXTENSIONS OF A VALUATION TO A TRANSCENDENTAL EXTENSION}

\subsection{The case of a monogenous transcendental extension}

LEMMA 1. Let K be a field, v a valuation on K, Γ its order group, Γ' a totally ordered

group containing \(\Gamma\) and \(\xi\) an element of \(\Gamma'\) . There exists a unique valuation \(w\) on \(\mathbf{K}(\mathbf{X})\) such that, for \(\mathbf{P} = \sum_{j} a_j \mathbf{X}^j (a_j \in \mathbf{K})\) , \(w(\mathbf{P}) = \inf_j(v(a_j) + j\xi)\) .

By Proposition 4 of § 3, no. 2, it suffices to show that the formula

\[
w \left(\sum_ {j} a _ {j} \mathbf {X} ^ {j}\right) = \inf _ {j} (v (a _ {j}) + j \xi)\tag{1}
\]

defines a valuation on the ring K{[}X{]}. As

\[
v \left(a _ {j} + b _ {j}\right) + j \xi \geqslant \inf (v \left(a _ {j}\right), v \left(b _ {j}\right)) + j \xi = \inf (v \left(a _ {j}\right) + j \xi , v \left(b _ {j}\right) + j \xi),
\]

it follows that

\[
w (\mathbf {P} + \mathbf {Q}) \geqslant \inf (w (\mathbf {P}), w (\mathbf {Q}))\tag{2}
\]

for \(\mathbf{P},\mathbf{Q}\) in \(\mathbf{K}[\mathbf{X}]\) , equality holding if \(w(\mathbf{P})\neq w(\mathbf{Q})\) . We show that

\[
w (\mathrm{PQ}) = w (\mathrm{P}) + w (\mathrm{Q})\tag{3}
\]

for \(\mathbf{P} = \sum_{j} a_{j} \mathbf{X}^{j}\) and \(\mathbf{Q} = \sum_{j} b_{j} \mathbf{X}^{j}\) . Let \(i\) (resp. \(k\) ) be the least of the integers \(j\) such that \(v(a_{j}) + j\xi\) (resp. \(v(b_{j}) + j\xi\) ) attains its minimum; let \(\alpha\) (resp. \(\beta\) ) denote this minimum. For \(j, j'\) in \(\mathbf{N}\) ,

\[
w (a _ {j} b _ {j ^ {\prime}} \mathbf {X} ^ {j + j ^ {\prime}}) = v (a _ {j}) + j \xi + v (b _ {j ^ {\prime}}) + j ^ {\prime} \xi \geqslant \alpha + \beta ,
\]

whence \(w(\mathrm{PQ}) \geqslant \alpha + \beta\) by (2). Consider now the term \(c\mathbf{X}^{i + k}\) of degree \(i + k\) in PQ; then \(c = \sum_{n \in \mathbf{Z}} a_{i + n} b_{k - n}\) ; by the choice of \(i\) and \(k\) , the element

\[
w (a _ {i + n} b _ {k - n} \mathbf {X} ^ {i + k}) = v (a _ {i + n}) + (i + n) \xi + v (b _ {k - n}) + (k - n) \xi
\]

takes its minimum value \(\alpha + \beta\) once and once only with \(n = 0\) ; hence \(w(c\mathbf{X}^{i + k}) = \alpha + \beta\) , whence, by (1),

\[
w (\mathrm{PQ}) = \alpha + \beta = w (\mathrm{P}) + w (\mathrm{Q}).
\]

PROPOSITION 1. Let K be a field, v a valuation on K, \(\Gamma\) its order group, \(\Gamma'\) a totally ordered group containing \(\Gamma\) and \(\xi\) an element of \(\Gamma'\) such that the relations \(n\xi\in\Gamma\) , \(n\in Z\) imply n=0. Then there exists a unique valuation w on K(X) with values in \(\Gamma'\) and extending v such that \(w(\mathbf{X})=\xi\) . The residue field of w is equal to that of v and its order group is the subgroup \(\Gamma+Z\xi\) of \(\Gamma'\) .

We show first the uniqueness of \(w\) . Let \(\mathbf{P} = \sum_{j} a_j \mathbf{X}^j\) be an element of \(\mathbf{K}[\mathbf{X}]\) . Then \(w(a_j \mathbf{X}^j) = v(a_j) + j\xi\) , which shows that the monomials \(a_j \mathbf{X}^j\) such that \(a_j \neq 0\) have distinct values for \(w\) . It follows that \(w(\mathbf{P}) = \inf_j (v(a_j) + j\xi)\) , which shows both the uniqueness of \(w\) on \(\mathbf{K}[\mathbf{X}]\) (hence also on \(\mathbf{K}(\mathbf{X})\) ) and the fact that the order group of \(w\) is \(\Gamma + \mathbf{Z}\xi\) . It is further seen that, if \(\mathbf{P} \neq 0\) , we may write \(\mathbf{P} = a\mathbf{X}^n (1 + u)\) , where \(a\in \mathbf{K}^*\) , \(n\in \mathbf{N}\) , \(u\in \mathbf{K}(\mathbf{X})\) and \(w(u) > 0\) ; every element \(\mathbf{R}\neq 0\) of \(\mathbf{K}(\mathbf{X})\) can therefore be written in the form

\[
\mathrm{R} = b \mathrm{X} ^ {n} (1 + u ^ {\prime}),
\]

where \(b \in \mathbf{K}^*\) , \(n \in \mathbf{Z}\) , \(u' \in \mathbf{K}(\mathbf{X})\) and \(w(u') > 0\) ; then \(w(\mathbf{R}) = v(b) + n\xi\) , hence \(w(\mathbf{R}) = 0\) if and only if \(v(b) = 0\) and \(n = 0\) ; thus, when \(w(\mathbf{R}) = 0\) , \(\mathbf{R}\) and \(b\) are congruent modulo the ideal of \(w\) , which shows that the residue field of \(w\) is equal to that of \(v\) .

Finally the existence of w follows from Lemma 1.

PROPOSITION 2. Let K be a field, v a valuation on K, \(\Gamma\) its order group and k its residue field. There exists a unique valuation w on K(X) extending v such that \(w(X) = 0\) and the image t of X in the residue field \(k'\) of w is transcendental over k. The order group of w is equal to that of v and its residue field is \(k(t)\) .

To show the uniqueness of \(w\) , it will suffice for us to show that, if \(\mathbf{P} = \sum_{j} a_j \mathbf{X}^j\) is a non-zero element of \(\mathbf{K}[\mathbf{X}]\) , then

\[
w (\mathbf {P}) = \inf _ {j} (v (a _ {j})).
\]

We may divide P by an element of \(K^{*}\) and suppose that \(v(a_{j}) \geqslant 0\) for all j and that one of the \(v(a_{j})\) is zero. As \(w(\mathbf{X}) = 0\) , P then belongs to the ring of w; writing \(\bar{a}_{j}\) for the canonical image of \(a_{j}\) in k, the canonical image of P in the residue field \(k'\) is \(\sum_{j} \bar{a}_{j} t^{j}\) ; as t is transcendental over k and one of the \(\bar{a}_{j}\) is non-zero, this image is non-zero, whence

\[
w (\mathbf {P}) = 0 = \inf _ {j} (v (a _ {j})).
\]

We now show the existence of \(w\) . The formula \(w(\mathbf{P}) = \inf_{j}(v(a_j))\) (for \(\mathbf{P} = \sum_{j} a_j \mathbf{X}^j\) ) defines a valuation \(w\) on \(\mathbf{K}(\mathbf{X})\) , by virtue of Lemma 1, and \(w\) obviously has the same order group as \(v\) . Then \(w(\mathbf{X}) = 0\) . We show that the canonical image \(t\) of \(\mathbf{X}\) in the residue field \(k'\) of \(w\) is transcendental over \(k\) : if \(\sum_{j} \bar{a}_j t^j = 0\) , where \(\bar{a}_j \in k\) for all \(j\) , then, denoting by \(a_j\) a representative of \(\bar{a}_j\) in the ring of \(v\) , \(w\left(\sum_{j} a_j \mathbf{X}^j\right) > 0\) ; whence \(v(a_j) > 0\) for all \(j\) , hence \(\bar{a}_j = 0\) for all \(j\) . We show finally that \(k' = k(t)\) : every element \(\mathbf{R}\) of \(\mathbf{K}(\mathbf{X})\) may be written \(\mathbf{R} = c\left(\sum_{j} a_j \mathbf{X}^j\right) / \left(\sum_{j} b_j \mathbf{X}^j\right)\) , where \(c, a_j, b_j\) are in \(\mathbf{K}\) , \(v(a_j) \geqslant 0\) and \(v(b_j) \geqslant 0\) for all \(j\) , one of the \(v(a_j)\) and one of the \(v(b_j)\) being zero; then \(w(\mathbf{R}) \geqslant 0\) if and only if \(v(c) \geqslant 0\) ; denoting by \(f\) the canonical homomorphism of the ring of \(w\) onto \(k'\) ,

\[
f (\mathrm{R}) = f (c) \left(\sum_ {j} f (a _ {j}) t ^ {j}\right) / \left(\sum_ {j} f (b _ {j}) t ^ {j}\right),
\]

which proves our assertion.

Remark. It should not be thought that the two types of extensions of v to K(X) which we have just met are the only ones; there may exist a third type of extension, where \(\Gamma'/\Gamma\) is a torsion group and \(k'\) a (not necessarily finite) algebraic extension of k. This third type is not necessarily provided by the procedure described in Lemma 1 (cf.~§3, Exercise 1).

\subsection{The rational rank of commutative group}

DEFINITION 1. The rational rank of a commutative group G is the dimension of the vector Q-space \(G \otimes_{\mathbf{Z}} \mathbf{Q}\).

This dimension may also be defined as the least upper bound (finite or infinite) of the cardinals r such that there exist r elements of G which are linearly independent over Z (Algebra, Chapter II, § 7, no. 10, Proposition 26). The rational rank of G is zero if and only if G is a torsion group. For a subgroup of an additive group \(R^{n}\) , the notion of rational rank coincides with that defined in General Topology, Chapter VII, § 1.

In the rest of this paragraph, we shall denote by \(r(G)\) the rational rank of the commutative group \(G\) . If \(G'\) is a subgroup of \(G\) , then (since \(Q\) is a flat \(Z\) -module) we have the additive equation

\[
r (\mathbf {G}) = r \left(\mathbf {G} ^ {\prime}\right) + r \left(\mathbf {G} / \mathbf {G} ^ {\prime}\right).\tag{4}
\]

PROPOSITION 3. Let G be a totally ordered commutative group and H a subgroup of G. If \(h(\mathbf{G})\) and \(h(\mathbf{H})\) denote the heights of G and H (§ 4, no. 4), then the inequality

\[
h (\mathbf {G}) \leqslant h (\mathbf {H}) + r (\mathbf {G} / \mathbf {H})\tag{5}
\]

holds.

In fact, let \(G_{0} \subset G_{1} \subset \cdots \subset G_{n}\) be a strictly increasing sequence of isolated subgroups of G. It is necessary to establish the inequality

\[
n \leqslant h (\mathrm{H}) + r (\mathrm{G} / \mathrm{H}).\tag{6}
\]

It is obvious for n = 0. Suppose \(n \geqslant 1\) and let us argue by induction on n.~Applying the induction hypothesis to the group \(G_{n-1}\) and its subgroup \(H \cap G_{n-1}\) , we obtain

\[
n - 1 \leqslant h (\mathrm{H} \cap \mathrm{G} _ {n - 1}) + r (\mathrm{G} _ {n - 1} / (\mathrm{H} \cap \mathrm{G} _ {n - 1})).\tag{7}
\]

Then we distinguish two cases:

\begin{enumerate}
\def\labelenumi{(\alph{enumi})}
\tightlist
\item
  \(\mathbf{H} \cap \mathbf{G}_{n-1} = \mathbf{H}\) , in other words \(\mathbf{H} \subset \mathbf{G}_{n-1}\) . Inequality (7) may be written
\end{enumerate}

\[
n \leqslant h (\mathrm{H}) + r (\mathrm{G} _ {n - 1} / \mathrm{H}) + 1.\tag{8}
\]

Now \(\mathbf{G} / \mathbf{G}_{n - 1}\) is a totally ordered group not reduced to 0; it is therefore not a torsion group and \(r(\mathrm{G} / \mathrm{G}_{n - 1}) \geqslant 1\) . Whence by (4), \(r(\mathrm{G} / \mathrm{H}) \geqslant r(\mathrm{G}_{n - 1} / \mathrm{H}) + 1\) . Substituting in (8), we certainly obtain the desired inequality (6).

\begin{enumerate}
\def\labelenumi{(\alph{enumi})}
\setcounter{enumi}{1}
\tightlist
\item
  \(\mathbf{H} \cap \mathbf{G}_{n-1} \neq \mathbf{H}\) . As \(\mathbf{H} \cap \mathbf{G}_{n-1}\) is an isolated subgroup of \(\mathbf{H}\) , we conclude that \(h(\mathbf{H}) \geqslant h(\mathbf{H} \cap \mathbf{G}_{n-1}) + 1\) . On the other hand, obviously \(r(\mathbf{G}/\mathbf{H}) \geqslant r(\mathbf{G}_{n-1}/(\mathbf{H} \cap \mathbf{G}_{n-1}))\) . Substituting in (7), we again obtain (6).
\end{enumerate}

COROLLARY. For every totally ordered commutative group \(\mathbf{G}\) , \(h(\mathbf{G}) \leqslant r(\mathbf{G})\) .

We set \(\mathbf{H} = \{0\}\) in Proposition 3.

PROPOSITION 4. Let G be a totally ordered commutative group. Suppose that G is finitely generated and that \(h(\mathbf{G}) = r(\mathbf{G})\) . Then G is isomorphic to \(\mathbf{Z}^{r(\mathbf{G})}\) ordered lexicographically.

Let us write \(r = r(\mathbf{G}) = h(\mathbf{G})\) . If \(r = 0\) , then \(\mathbf{G} = \{0\}\) . If \(r = 1\) , the structure of finitely generated commutative groups shows that there is an isomorphism \(j\) of \(\mathbf{G}\) onto \(\mathbf{Z}\) (Algebra, Chapter VII, § 4, no. 6, Theorem 3). Now \(\mathbf{Z}\) has only two total orderings compatible with its group structure, namely the usual ordering and its opposite. Hence \(j\) or \(-j\) is an isomorphism of the ordered group \(\mathbf{G}\) onto \(\mathbf{Z}\) with the usual ordering.

Suppose now that \(r \geqslant 2\) and let us argue by induction on \(r\) . Let \(H\) be an isolated subgroup of \(G\) of height \(r - 1\) . Then \(r(H) + r(G/H) = r\) (formula (4)), \(r(H) \geqslant h(H) = r - 1\) and \(r(G/H) \geqslant h(G/H) = 1\) (Corollary to Proposition 3), whence \(r(H) = r - 1\) and \(r(G/H) = 1\) . The induction hypothesis shows that \(H\) is isomorphic to \(Z^{r-1}\) ordered lexicographically and the case \(r = 1\) shows that \(G/H\) is isomorphic to \(Z\) . As \(Z\) is a free \(Z\) -module, \(H\) is a direct factor of \(G\) (Algebra, Chapter II, § 1, no. 11, Proposition 21). The following lemma then shows that \(G\) is isomorphic (not canonically) to the lexicographical product \(H \times (G/H)\) , which completes the proof.

LEMMA 2. Let H be an isolated subgroup of a totally ordered commutative group G. If H is a direct factor of G, the ordered group G is isomorphic to the group (G/H) × H ordered lexicographically.

Let \(j\) be an isomorphism of \((\mathrm{G} / \mathrm{H}) \times \mathrm{H}\) onto \(\mathbf{G}\) such that \(j(0, x) = x\) for all \(x \in \mathrm{H}\) and \(j(y, x)\) belongs to the coset \(y\) of \(\mathrm{H}\) . As \((\mathrm{G} / \mathrm{H}) \times \mathrm{H}\) is totally ordered, it amounts to showing that \(j\) is increasing (Set Theory, Chapter III, § 1, no. 12, Proposition 11). Let \((y, x)\) be an element \(\geqslant 0\) of \((\mathrm{G} / \mathrm{H}) \times \mathrm{H}\) ordered lexicographically. If \(y > 0\) , the coset of \(\mathrm{H}\) containing \(j(y, x)\) is an element \(> 0\) , whence \(j(y, x) > 0\) , for, otherwise, \(y \leqslant 0\) (§ 4, no. 2, Proposition 3). If \(y = 0\) and \(x \geqslant 0\) , then \(j(y, x) = x \geqslant 0\) . Hence \(j\) is certainly increasing.

\subsection{The case of any transcendental extension}

In this no. we shall use the following notation: K is a field, \(K'\) an extension of K, v a valuation on K, \(v'\) an extension of v to \(K'\) , \(\Gamma\) and k (resp. \(\Gamma'\) and \(k'\) ) the order group and residue field of \(v\) (resp. \(v'\) ). We shall write:

\[
d \left(\mathrm{K} ^ {\prime} / \mathrm{K}\right) = \dim . \mathrm{al} _ {\mathrm{K}} \mathrm{K} ^ {\prime} = \text { transcendence   degree   of } \mathrm{K} ^ {\prime} \text { over } \mathrm{K};
\]

\[
s (v ^ {\prime} / v) = \dim . \mathrm{al} _ {k} k ^ {\prime} = \text { transcendence   degree   of } k ^ {\prime} \text { over } k;
\]

\[
r (v ^ {\prime} / v) = r (\Gamma^ {\prime} / \Gamma) = \text { rational   rank   of } \Gamma^ {\prime} / \Gamma ,
\]

if the right-hand sides are finite; otherwise, we shall make the convention that \(d(\mathbf{K}' / \mathbf{K}) = +\infty\) (resp. \(s(v' / v) = +\infty\) , \(r(v' / v) = +\infty\) ).

THEOREM 1. Let \(x_1, \ldots, x_s\) be elements of the ring of \(v'\) whose canonical images \(\bar{x}_i\) in \(k'\) are algebraically independent over \(k\) and \(y_1, \ldots, y_r\) elements of \(K'\) such that the canonical images of the \(v'(y_j)\) in \(\Gamma'/\Gamma\) are linearly independent over \(Z\) . Then the \(r + s\) elements \(x_1, \ldots, x_s, y_1, \ldots, y_r\) of \(K'\) are algebraically independent over \(K\) ; the restriction of \(v'\) to \(K(x_1, \ldots, x_s, y_1, \ldots, y_r)\) admits \(k(\bar{x}_1, \ldots, \bar{x}_s)\) as residue field and

\[
\Gamma + \mathbf {Z} v ^ {\prime} (y _ {1}) + \dots + \mathbf {Z} v ^ {\prime} (y _ {r})
\]

as order group.

Our assertion is obvious if \(r + s = 0\) . We argue by induction on \(r + s\) . If \(r' \leqslant r\) , \(s' \leqslant s\) and \(r' + s' < r + s\) , the induction hypothesis shows that the hypotheses of Theorem 1 hold if \(K\) is replaced by \(K(x_1, \ldots, x_s, y_1, \ldots, y_{r'})\) and the families \((x_1, \ldots, x_s), (y_1, \ldots, y_r)\) by \((x_{s'+1}, \ldots, x_s), (y_{r'+1}, \ldots, y_r)\) . The problem is therefore reduced to one of the two following cases:

\begin{enumerate}
\def\labelenumi{(\alph{enumi})}
\item
  There is an element \(x\) in the ring of \(v'\) such that \(\bar{x}\) is transcendental over \(k\) ; then it is necessary to show that \(x\) is transcendental over \(\mathbf{K}\) and that the restriction of \(v'\) to \(\mathbf{K}(x)\) admits \(k(\bar{x})\) as residue field and \(\Gamma\) as order group.
\item
  There is an element \(y\) in \(\mathbf{K}'\) such that the relations \(nv'(y) \in \Gamma\) and \(n \in \mathbf{Z}\) imply \(n = 0\) ; it is necessary to show that \(y\) is transcendental over \(\mathbf{K}\) and that the restriction of \(v'\) to \(\mathbf{K}(y)\) admits \(k\) as residue field and \(\Gamma + \mathbf{Z}v'(y)\) as order group.
\end{enumerate}

Now Proposition 1 of §8, no. 1 shows that \(x\) (resp. \(y\) ) cannot be algebraic over \(K\) . The other assertions of (a) (resp. (b)) follow immediately by Proposition 2 (resp. Proposition 1) of no. 1.

COROLLARY 1. The inequality

\[
s (v ^ {\prime} / v) + r (v ^ {\prime} / v) \leqslant d (\mathbf {K} ^ {\prime} / \mathbf {K})\tag{9}
\]

holds.

Further, if \(K'\) is a finitely generated extension of K and equality holds in (9), then \(\Gamma'/\Gamma\) is a finitely generated Z-module and \(k'\) is a finitely generated extension of k.

Let \(r\) and \(s\) be natural numbers such that \(r \leqslant r(v'/v)\) and \(s \leqslant s(v'/v)\) ; we show that \(r + s \leqslant d(\mathbf{K}'/\mathbf{K})\) and this will prove (9). By hypothesis there exist elements \(x_1, \ldots, x_s, y_1, \ldots, y_r\) of \(\mathbf{K}'\) which satisfy the hypotheses of Theorem 1.

They are therefore algebraically independent over K, which shows the inequality \(r + s \leqslant d(\mathbf{K}'/\mathbf{K})\) .

If \(\mathbf{K}'\) is a finitely generated extension of \(\mathbf{K}\) , \(d(\mathbf{K}' / \mathbf{K})\) is finite, hence \(s(v' / v)\) and \(r(v' / v)\) are also finite; we denote them by \(s\) and \(r\) . There exist elements \(x_{1},\ldots ,x_{s},\) \(y_{1},\ldots ,y_{r}\) of \(\mathbf{K}'\) which satisfy the hypotheses of Theorem 1. If \(r + s = d(\mathbf{K}' / \mathbf{K})\) , these elements form a transcendence basis of \(\mathbf{K}'\) over \(\mathbf{K}\) and \(\mathbf{K}'\) is therefore a finite algebraic extension of \(\mathbf{K}'' = \mathbf{K}(x_1,\dots,y_r)\) . Let \(\Gamma''\) and \(k''\) be the order group and residue field of the restriction of \(v'\) to \(\mathbf{K}''\) . By Theorem 1, \(\Gamma'' / \Gamma\) is a finitely generated \(\mathbf{Z}\) -module and \(k''\) is a finitely generated pure extension of \(k\) . On the other hand, as \(\mathbf{K}'\) is a finite algebraic extension of \(\mathbf{K}''\) , \(\Gamma' / \Gamma\) is a finite group and \(k'\) is a finite algebraic extension of \(k''\) (§ 8, no. 1, Lemma 2). This proves the corollary.

COROLLARY 2. Let \(h\) and \(h'\) be the heights of \(v\) and \(v'\) . Then

\[
s (v ^ {\prime} / v) + h ^ {\prime} \leqslant d (\mathrm{K} ^ {\prime} / \mathrm{K}) + h.\tag{10}
\]

By Proposition 3, \(h' \leqslant r(v'/v) + h\) .

COROLLARY 3. Suppose that \(\mathbf{K}'\) is a finitely generated extension of \(\mathbf{K}\) , that \(\Gamma\) is isomorphic to \(\mathbf{Z}^h\) (ordered lexicographically) and there is equality in formula (10). Then \(\Gamma'\) is isomorphic to \(\mathbf{Z}^{h'}\) (ordered lexicographically) and \(k'\) is a finitely generated extension of \(k\) .

If there is equality in (10), there is equality in (9), whence the fact that \(k'\) is a finitely generated extension of k and that \(\Gamma'\) is a finitely generated Z-module. Further, comparing (9) and (10), it is seen that \(h' - h = r(\Gamma'/\Gamma)\) , whence \(h' = r(\Gamma')\) and Proposition 4 (no. 2) then shows that \(\Gamma'\) is isomorphic to \(Z^{h'}\) ordered lexicographically.

COROLLARY 4. Suppose that \(v\) is improper (in which case \(k = K\) ). Then

\[
h (\Gamma^ {\prime}) + d (k ^ {\prime} / \mathbf {K}) \leqslant r (\Gamma^ {\prime}) + \mathbf {D} (k ^ {\prime} / \mathbf {K}) \leqslant d (\mathbf {K} ^ {\prime} / \mathbf {K}).\tag{11}
\]

If, in particular, \(v'\) is of height 1, then

\[
d (k ^ {\prime} / \mathbf {K}) \leqslant d (\mathbf {K} ^ {\prime} / \mathbf {K}) - 1;\tag{12}
\]

further, if \(\mathbf{K}'\) is a finitely generated extension of \(\mathbf{K}\) and there is equality in (12), then \(v'\) is a discrete valuation and \(k'\) is a finitely generated extension of \(\mathbf{K}\) .

It is a series of special cases of Corollaries 1, 2, 3.


\exercisesection{1}

\begin{enumerate}
\def\labelenumi{\arabic{enumi}.}
\tightlist
\item
  Let K be a field; for every subring A of K let L(A) denote the set of local rings \(A_{p}\) of A where p runs through the prime ideals of A (these local rings being identified with subrings of K).
\end{enumerate}

\begin{enumerate}
\def\labelenumi{(\alph{enumi})}
\item
  For a local subring M of K to dominate a ring \(A_{p} \in L(A)\) , it is necessary and sufficient that \(A \subset M\) ; the local ring \(A_{p}\) dominated by M is then unique and corresponds to the prime ideal \(p = m(M) \cap A\) .
\item
  Let M, N be two local subrings of K and P the subring of K generated by \(M \cup N\) . Show that the following conditions are equivalent:
\end{enumerate}

\((\alpha)\) There exists a prime ideal \(\mathfrak{p}\) of \(\mathbf{P}\) such that \(\mathfrak{m}(\mathbf{M}) = \mathfrak{p} \cap \mathbf{M}\) , \(\mathfrak{m}(\mathbf{N}) = \mathfrak{p} \cap \mathbf{N}\) .

(β) The ideal \(\mathfrak{a}\) generated in \(\mathbf{P}\) by \(\mathfrak{m}(\mathbf{M}) \cup \mathfrak{m}(\mathbf{N})\) is distinct from \(\mathbf{P}\) .

(γ) There exists a local subring Q of K dominating both M and N.

If these conditions are satisfied, M and N are called related.

\begin{enumerate}
\def\labelenumi{(\alph{enumi})}
\setcounter{enumi}{2}
\tightlist
\item
  Let A, B be two subrings of K and C the subring of K generated by A ∪ B. Show that the following conditions are equivalent:
\end{enumerate}

\((\alpha)\) For every local ring \(\mathbf{Q} \subset \mathbf{K}\) containing A and B, \(\mathbf{A}_{\mathfrak{p}} = \mathbf{B}_{\mathfrak{q}}\) , where \(\mathfrak{p} = \mathfrak{m}(\mathbf{Q}) \cap \mathbf{A}\) , \(\mathfrak{q} = \mathfrak{m}(\mathbf{Q}) \cap \mathbf{B}\) .

(β) For every prime ideal r of C, \(A_{p} = B_{q}\) , where \(p = r \cap A\) , \(q = r \cap B\) .

(γ) If M ∈ L(A) and N ∈ L(B) are related, they are identical.

(δ) L(A) ∩ L(B) = L(C).

(Use (a) and (b)).

\begin{enumerate}
\def\labelenumi{\arabic{enumi}.}
\setcounter{enumi}{1}
\item
  For an integral domain A to be a valuation ring, it is necessary and sufficient that every ideal of A be irreducible (Algebra, Chapter II, § 2, Exercise 16).
\item
  Show that every valuation ring is coherent (Chapter I, § 2, Exercise 12).
\item
  Let K be a field and F a set of valuation rings of K totally ordered by inclusion. Show that the intersection of the rings belonging to F is a valuation ring of K.
\item
  Let K be a field, A a integrally closed subdomain of K with K as a field of fractions and \((\mathbf{A}_{\alpha})\) a family of valuation rings of K whose intersection is A. Then if L is an extension of K, the integral closure of A in L is the intersection of the valuation rings of L which dominate one of the \(A_{\alpha}\) .
\end{enumerate}

¶ 6. Let R be a ring and A a subring of R such that S = R - A is non-empty and the product of two elements of S belongs to S.

\begin{enumerate}
\def\labelenumi{(\alph{enumi})}
\item
  Let \(a, a'\) be two elements of \(\mathbf{A}\) and \(s, s'\) two elements of \(S\) such that \(sa \in \mathbf{A}\) and \(s'a' \in \mathbf{A}\) . Show that one of the two elements \(sa', s'a\) belongs to \(\mathbf{A}\) . Deduce that the set \(p_1\) of elements \(a \in \mathbf{A}\) for which there exists \(s \in S\) such that \(sa \in \mathbf{A}\) is a prime ideal of \(\mathbf{A}\) .
\item
  Show that the set \(\mathfrak{p}_0\) of \(a \in \mathbf{A}\) such that \(sa \in \mathbf{A}\) for all \(s \in \mathbf{S}\) is a prime ideal of \(\mathbf{R}\) and \(\mathbf{A}\) ; it is the greatest ideal of \(\mathbf{R}\) contained in \(\mathbf{A}\) and \(\mathfrak{p}_0 \subset \mathfrak{p}_1\) .
\item
  The set \(\mathfrak{n}\) of elements \(x \in \mathbb{R}\) such that there exists \(s \in \mathbb{S}\) for which \(sx = 0\) is both an ideal of \(\mathbb{R}\) and an ideal of \(\mathbb{A}\) and \(\mathfrak{n} \subset \mathfrak{p}_0\) .
\item
  The ring \(\mathbf{A}\) is integrally closed in \(\mathbb{R}\) .
\item
  If \(\mathbf{R}\) is a field, \(\mathbf{A}\) is a valuation ring of \(\mathbf{R}\) , \(p_0 = (0)\) and \(p_1\) is the maximal ideal of \(\mathbf{A}\) .
\end{enumerate}

¶ 7. Let R be a ring. An ordered pair (A, p) consisting of a subring A of R and a prime ideal p of A is called maximal in R if there exists no ordered pair (A', p') consisting of a subring A' of R and a prime ideal p' of A' such that A ⊂ A', A' ≠ A and p' ∩ A = p.

\begin{enumerate}
\def\labelenumi{(\alph{enumi})}
\item
  For every subring B of R and every prime ideal q of B there exists an ordered pair \((A, p)\) , maximal in R, such that \(A \supset B\) and \(q = p \cap B\) .
\item
  For an ordered pair (A, p) to be maximal in R, it is necessary and sufficient that, for all \(s \in R - A\) , there exist a finite family of elements \(c_{i} \in p\) ( \(1 \leqslant i \leqslant n\) ) such that the element \(b = c_{1}s + c_{2}s^{2} + \cdots + c_{n}s^{n}\) belongs to A - p (use Chapter II, §2, no. 5, Corollary 3 to Proposition 11).
\item
  Show that, if the ordered pair \((\mathbf{A},\mathfrak{p})\) is maximal in \(\mathbf{R}\) , the product of two elements of \(\mathbf{R} - \mathbf{A}\) belongs to \(\mathbf{R} - \mathbf{A}\) . (If \(s,s'\) are in \(\mathbf{R} - \mathbf{A}\) and \(ss' = a\in \mathbf{A}\) , consider the least integers \(n,m\) such that there exists \(c_{i}\in \mathfrak{p}\) ( \(1\leqslant i\leqslant n\) ) for which \(b = \sum_{i}c_{i}s^{i}\in \mathbf{A} - \mathfrak{p}\) and \(c_{j}'\in \mathfrak{p}\) ( \(1\leqslant j\leqslant m\) ) for which
\end{enumerate}

\[
b ^ {\prime} = \sum_ {j} c _ {j} ^ {\prime} s ^ {\prime j} \in A - p;
\]

supposing for example \(n \geqslant m\) , consider the product \(bb' \in A - p\) and obtain a contradiction from the definition of the integer \(n\) .)

\begin{enumerate}
\def\labelenumi{(\alph{enumi})}
\setcounter{enumi}{3}
\item
  Under the hypotheses of (c), we write \(S = R - A\) ; show that the ideal \(p_1\) of \(A\) consisting of the \(a \in A\) such that there exists \(s \in S\) for which \(sa \in A\) is contained in \(p\) (use (b)). A fortiori the ideals \(p_0\) of \(A\) consisting of the \(a \in A\) such that \(sa \in A\) for all \(s \in S\) is contained in \(p\) (Exercise 6 (b)). We write \(A' = A / p_0\) , \(p' = p / p_0\) and \(R' = R / p_0\) ; the ordered pair \((A', p')\) is then maximal in the integral domain \(\mathbf{R}'\) and, for all non-zero \(a' \in \mathbf{A}'\) , there exists \(s' \in \mathbf{S}' = \mathbf{R}' - \mathbf{A}'\) such that \(s'a' \in \mathbf{S}'\) .
\item
  Let \((\mathbf{A},\mathfrak{p})\) be a maximal ordered pair in the integral domain \(\mathbf{R}\) such that, for all non-zero \(a\in \mathbf{A}\) , there exists \(s\in S = R - A\) such that \(sa\in S\) . We write \(\mathrm{S}_0 = \mathrm{S}\cup \{1\}\) and denote by \(\mathbb{R}_0\) the ring of fractions \(\mathrm{S}_0^{-1}\mathrm{R}\) ; the canonical homomorphism \(\mathbb{R}\to \mathbb{R}_0\) is injective and therefore identifies \(\mathbb{R}\) with a subring of \(\mathbb{R}_0\) . Show that \(\mathbb{R}_0\) is a field identified with the field of fractions of \(\mathbb{R}\) . Let \((\mathbf{B},\mathfrak{q})\) be a maximal ordered pair in \(\mathbb{R}_0\) such that \(\mathbf{B}\supset \mathbf{A}\) and \(\mathfrak{q}\cap \mathbf{A} = \mathfrak{p}\) . Show that \(\mathbf{B}\cap \mathbf{R} = \mathbf{A}\) ; in other words, \(\mathbf{A}\) is the intersection of \(\mathbf{R}\) and a valuation ring of \(\mathbb{R}_0\) and \(\mathfrak{p}\) is the intersection of \(\mathbf{A}\) and the maximal ideal of this valuation ring.
\item
\end{enumerate}

¶ 8. Given a ring A, a polynomial P(X \(_{1}\) , \ldots, X \(_{n}\) ) with coefficients in A is called dominated if it is of the form X \(^{\alpha}\) + ∑ \(_{\beta} c_{\beta}\) X \(^{\beta}\) , where β \textless{} α in the product ordered set N \(^{n}\) for all β such that c \(_{\beta}\) ≠ 0. In a ring R, a subring A is called paravaluative if, for every dominated polynomial P ∈ A{[}X \(_{1}\) , \ldots, X \(_{n}\) {]} (any n), P(s \(_{1}\) , \ldots, s \(_{n}\) ) ≠ 0 for all elements s \(_{i}\) ∈ R − A (1 ≤ i ≤ n).

\begin{enumerate}
\def\labelenumi{(\alph{enumi})}
\item
  Show that, if \(\mathbf{A}\) is paravaluative in \(\mathbb{R}\) , the product of two elements of \(\mathbf{S} = \mathbf{R} - \mathbf{A}\) belongs to \(\mathbf{S}\) .
\item
  Let A be a subring of a ring R such that the product of two elements of S = R - A belongs to S; let \(p_{0}\) be the prime ideal of A and R consisting of the \(a \in A\) such that \(sa \in A\) for all \(s \in S\) (Exercise 6 (b)). For A to be paravaluative in R, it is necessary and sufficient that \(A/p_{0}\) be paravaluative in \(R/p_{0}\) .
\item
  Let R be a ring and \(h: \mathbf{R} \to \mathbf{R}'\) a ring homomorphism. If \(\mathbf{A}'\) is a paravaluative subring of \(\mathbf{R}'\) , \(\mathbf{A} = \bar{h}^{-1}(\mathbf{A}')\) is paravaluative in R. Deduce that, if \((\mathbf{A}, p)\) is a maximal ordered pair in R (Exercise 7), A is paravaluative in R (use (b) and Exercise 7 (d) and (e)).
\item
  Let \(\mathbf{R}'\) be a ring and \(\mathbf{R}\) a subring of \(\mathbf{R}'\) . If \(\mathbf{A}\) is a paravaluative subring of \(\mathbf{R}\) , show that there exists a paravaluative subring \(\mathbf{A}'\) of \(\mathbf{R}'\) such that \(\mathbf{A}' \cap \mathbf{R} = \mathbf{A}\) . (If \(S = R - A\) and \(S_0 = S \cup \{1\}\) , consider the ring \(T' = S_0^{-1}R'\) and the canonical homomorphism \(h: R' \to T'\) . Let \(\mathbf{B}\) be the subring of \(T'\) generated by \(h(A) \cup (h(S))^{-1}\) and \(\mathfrak{b}\) the ideal of \(\mathbf{B}\) generated by \((h(S))^{-1}\) . Show that \(\mathfrak{b} \neq \mathbf{B}\) ; consider a maximal ordered pair \((\mathbf{A}^{\prime \prime}, \mathfrak{p}^{\prime \prime})\) in \(\mathbf{T}'\) such that \(\mathbf{B} \subset \mathbf{A}''\) and \(\mathfrak{b} \subset \mathfrak{p}''\) and take \(\mathbf{A}' = \overline{h}^{-1}(\mathbf{A}'')\) .
\end{enumerate}

Conclude from this that, if R is an integral domain, the paravaluative sub-rings of R are the intersections of R with the valuation rings of the field of fractions of R.

¶ 9. (a) Let R be a ring, A a subring of R, x an element of R and S the set of \(x^{k}\) ( \(k \geqslant 0\) ) in R. Show that, for x to be integral over A, it is necessary and sufficient that \(x/1 \in A[1/x]\) in the ring \(S^{-1}R\) . If x is not integral over A, show that there exists a maximal ideal m of \(A[1/x]\) containing 1/x and that the inverse image of m under the canonical homomorphism \(A \to A[1/x]\) is a maximal ideal of A.

\begin{enumerate}
\def\labelenumi{(\alph{enumi})}
\setcounter{enumi}{1}
\tightlist
\item
  Show that the integral closure A' of A in R is the intersection of the paravaluative subrings of R (Exercise 8) which contain A. (To see that A' is contained in each of these rings, use Exercise 8 (a) and Exercise 6 (d); to see that, if \(x \in R\) is not integral over A, there exists a subring B of R which is paravaluative in R and such that \(x \notin B\) , use (a), Exercise 7 (a) and Exercise 8 (c), arguing as in Theorem 3 of no. 3.)
\end{enumerate}

\exercisesection{2}

\begin{enumerate}
\def\labelenumi{\arabic{enumi}.}
\item
  Let \(\Omega\) be an extension of a field K, E and F two extensions of K contained in \(\Omega\) and linearly disjoint over K and L an extension of K contained in E. Let \(f\) be a place of E with values in L which reduces to the identity automorphism on K. Show that there exists a unique place of K(E \(\cup\) F) with values in K(L \(\cup\) F), extending \(f\) and reducing to the identity automorphism on F. (If A is the ring of \(f\) , note that \(f\) can be extended uniquely to a homomorphism \(g\) from F{[}A{]} to K(L \(\cup\) F) reducing to the identity on F and that, if m is the ideal of \(f\) , every element \(z \in\) F{[}A{]} such that \(g(z) = 0\) can be written as \(xu\) , where \(x \in m\) , \(u \in\) F{[}A{]} and \(g(u) \neq 0\) . Conclude by noting that K(E \(\cup\) F) is the field of fractions of F{[}A{]}.)
\item
  Let K, K' be two extensions of a field k and f and \(f'\) places respectively of K and \(K'\) with values in the same algebraically closed field L. Suppose that the restrictions of f and \(f'\) to k coincide. Show that there exist a composite extension (F, i, i') of K and \(K'\) (Algebra, Chapter VIII, § 8) and a place g of F with values in L such that \(f(x) = g(i(x))\) for \(x \in \mathbf{K}\) and \(f'(x) = g(i'(x'))\) for \(x' \in \mathbf{K}'\) . (Let V, \(V'\) be the rings of f and \(f'\) , A their common intersection with k and h: \(V \otimes_{A} V' \to L\) the homomorphism such that \(h(a \otimes a') = f(a)f'(a')\) for \(a \in V\) , \(a' \in V'\) . Using the fact that V and \(V'\) are flat A-modules (§ 3, no. 5, Lemma 1), prove that, if \(a \neq 0\) , \(a' \neq 0\) , \(a \otimes a' = (a \otimes 1)(1 \otimes a')\) is not a divisor of zero; if S is the multiplicative subset of B = V \(\otimes_{A} V'\) consisting of these elements,
\end{enumerate}

\(S^{-1}B\) contains subfields \(K_{1}, K_{1}'\) respectively isomorphic to K and \(K'\) and is the ring generated by \(K_{1} \cup K_{1}'\) . Show that, if q is the kernel of h, there exists a prime ideal p in B such that \(p \subset q\) and \(p \cap S = \varnothing\) and show that F may be taken as the field of fractions of \(S^{-1}B/S^{-1}p\) .

\begin{enumerate}
\def\labelenumi{\arabic{enumi}.}
\setcounter{enumi}{2}
\tightlist
\item
  \begin{enumerate}
  \def\labelenumii{(\alph{enumii})}
  \tightlist
  \item
    Let \(\mathbf{K}\) be a field and \(f\) a place of \(\mathbf{K}\) with values in an orderable field \(\mathbf{L}\) (Algebra, Chapter VI, §2, Exercise 8); show that \(\mathbf{K}\) is orderable. (Note that, for a finite family of elements \((a_i)_{1\leq i\leq n}\) of \(\mathbf{K}\) , there is always an index \(j\) such that \(f(a_i / a_j)\neq \infty\) for all \(i\) .)
  \end{enumerate}
\end{enumerate}

If \((x_{i})_{1\leqslant i\leqslant n}\) is a sequence of elements of K such that \(f\left(\sum_{i = 1}^{n}x_{i}^{2}\right)\neq \infty\) , show that \(f(x_{i})\neq \infty\) for all \(i\)

\begin{enumerate}
\def\labelenumi{(\alph{enumi})}
\setcounter{enumi}{1}
\item
  Let K be an ordered field and G a subgroup of K; show that the ring \(F(G)\) of elements of K which are not infinitely great with respect to G (loc. cit., Exercise 11) is a valuation ring of K and the ideal \(I(G)\) of elements of \(F(G)\) infinitesimally small with respect to G is the maximal ideal of this ring; hence there is on K a place with values in \(k(G) = F(G)/I(G)\) , which reduces to the identity on G; this is called the canonical place of K with respect to G.
\item
  Show that, if there exists no extension of G contained in K, distinct from G and comparable with G, \(k(G)\) is algebraic over G (note that, if \(t \in K\) does not belong to G, the field \(G(t)\) contains an element \(u \neq 0\) infinitesimally small with respect to G and that \(G(t)\) is algebraic over \(G(u)\) ).
\end{enumerate}

¶ 4. An ordered field K is called Euclidean if, for all \(x \geqslant 0\) in K, there exists \(y \in \mathbf{K}\) such that \(x = y^2\) .

\begin{enumerate}
\def\labelenumi{(\alph{enumi})}
\item
  Let K be a Euclidean ordered field and G a subfield of K such that there exists no extension of G contained in K, distinct from G and comparable with G. Show that, if f is a place of K with values in an ordered field L, which is an algebraic extension of G, and reducing to the identity on G, then f is equivalent to the canonical place of K with respect to G. (f may be considered to take its values in a maximal ordered field N which is an algebraic extension of L. Observe that G is Euclidean and that, if \(x \in G\) , \(y \in K\) and x \textless{} y, then \(f(y) = \infty\) or \(f(y) > x\) . If A and p are the ring and ideal of the place f, show successively that \(p \subset I(G)\) , \(F(G) \subset A\) , \(I(G) \subset p\) and \(A \subset F(G)\) , noting that N is comparable with G.)
\item
  Let K be a Euclidean ordered field, f a place of K with values in a maximal ordered field L and A and p the ring and ideal of f.~Let G be a maximal subfield among the subfields E of K such that \(E \cap p = 0\) . Show that \(G \subset A\) and that G is algebraically closed in K and therefore Euclidean. Prove that there exists no extension of G contained in K, distinct from G and comparable with G. Moreover, the subfield of L generated by \(f(A)\) is algebraic over \(f(G)\) and f is therefore equivalent to the canonical place of K with respect to G.
\item
  Let K be a maximal ordered field, L a maximal ordered subfield of K, distinct from K and such that K is comparable with L (Algebra, Chapter VI, § 2, Exercise 15). Show that there exists no place f of K with values in L reducing to the identity automorphism on L.
\end{enumerate}

¶ 5. Let K be a field and f a place of K with values in a maximal ordered field L.

\begin{enumerate}
\def\labelenumi{(\alph{enumi})}
\item
  For all \(\alpha \in \mathbf{K}\) show that there exists an extension of \(f\) , either to \(\mathbf{K}(\sqrt{\alpha})\) or to \(\mathbf{K}(\sqrt{-\alpha})\) , which is a place with values in \(\mathbf{L}\) . (Consider two cases according to whether there exists \(a \in \mathbf{K}\) such that \(f(a^2\alpha)\) is neither 0 nor \(\infty\) or \(f(x^2\alpha)\) is equal to 0 or \(\infty\) for all \(x \in \mathbf{K}\) .)
\item
  Deduce from (a) that there exist an extension E of K which is a maximal ordered field and an extension of f to E which is a place of E with values in L. (Reduce it to the case where K is Euclidean and use Exercise 4 (a).)
\end{enumerate}

\begin{enumerate}
\def\labelenumi{\arabic{enumi}.}
\setcounter{enumi}{5}
\tightlist
\item
  Let A be an integrally closed domain, K its field of fractions, p a prime ideal of A and k the field of fractions of A/p.~Show that, if \(A_{p}\) is not a valuation ring, there exists a valuation ring V of K which dominates \(A_{p}\) and whose residue field is a transcendental extension of k (use Chapter V, §2, Exercise 14(b)).
\end{enumerate}

\exercisesection{3}

¶ 1. (a) Let L be a field and T an indeterminate; in the field of formal power series L((T)), consider the series

\[
s = c _ {0} + c _ {1} \mathrm{T} ^ {1!} + c _ {2} \mathrm{T} ^ {2!} + \dots + c _ {n} \mathrm{T} ^ {n!} + \dots
\]

where the \(c_{i} \in L\) are all \(\neq 0\) . Show that s is transcendental over the subfield \(\mathrm{L}(\mathrm{T})\) of \(\mathrm{L}((\mathrm{T}))\) . (Argue by reductio ad absurdum by supposing that there exists a relation \(a_{0}(\mathrm{T}) + a_{1}(\mathrm{T})s + \cdots + a_{q}(\mathrm{T})s^{q} = 0\) where \(a_{i}(\mathrm{T})\) are polynomials in \(\mathrm{L}(\mathrm{T})\) such that \(a_{q} \neq 0\) ; construct the equation satisfied by

\[
s ^ {\prime} = s - c _ {0} - c _ {1} \mathrm{T} ^ {1!} - \dots - c _ {n - 1} \mathrm{T} ^ {(n - 1)!}
\]

and show that for \(n\) sufficiently large its constant term is a polynomial of degree \(< n!\) , whence our contradiction.)

\begin{enumerate}
\def\labelenumi{(\alph{enumi})}
\setcounter{enumi}{1}
\tightlist
\item
  Let \(k\) be a field, \(\mathbf{X}, \mathbf{Y}, \mathbf{Z}\) , \(\mathrm{T}\) four indeterminates, \(\mathbf{K} = k(\mathbf{X}, \mathbf{Y}, \mathbf{Z})\) and \(\mathbf{E}\) an algebraic closure of \(k(\mathbf{X})\) ; in \(\mathbf{E}\) , for all \(n\) , let \(x_{n}\) denote an element such that \(x_{n}^{n} = \mathbf{X}\) . In the field of formal power series \(\mathrm{E}((\mathrm{T}))\) , consider the element
\end{enumerate}

\[
s = x _ {1} \mathrm{T} + x _ {2} \mathrm{T} ^ {2!} + \dots + x _ {n} \mathrm{T} ^ {n!} + \dots
\]

The elements \(s\) and \(\mathbf{T}\) are algebraically independent over \(\mathbf{E}\) by (a). Consider the homomorphism \(f\colon \mathbf{K} \to \mathbf{E}((\mathbf{T}))\) such that \(f(\mathbf{X}) = \mathbf{X}, f(\mathbf{Y}) = \mathbf{T}, f(\mathbf{Z}) = s\) . If \(v\) is the discrete valuation on \(\mathbf{E}((\mathbf{T}))\) defined by the order of formal power series (no. 3, Example 3), \(v \circ f\) is a discrete valuation on \(\mathbf{K}\) ; show that the residue field of \(v \circ f\) is the subfield \(k(x_1, x_2, \ldots, x_n, \ldots)\) of \(\mathbf{E}\) and is therefore infinite over \(k\) .

\begin{enumerate}
\def\labelenumi{\arabic{enumi}.}
\setcounter{enumi}{1}
\tightlist
\item
  Let \(\Gamma\) be a totally ordered group and \(k\) a field.
\end{enumerate}

\begin{enumerate}
\def\labelenumi{(\alph{enumi})}
\item
  Let A, B be two subsets of \(\Gamma\) well ordered by the induced ordering; show that the set \(A + B\) is well ordered and that, for all \(\gamma \in A + B\) , the set of ordered pairs \((\alpha, \beta) \in A \times B\) such that \(\alpha + \beta = \gamma\) is finite. (To prove the first assertion, argue by reductio ad absurdum by showing that otherwise it would be possible to define a strictly decreasing sequence \((\gamma_n)\) of elements of \(A + B\) such that \(\gamma_n = \alpha_n + \beta_n\) , where \(\alpha_n \in A\) , \(\beta_n \in B\) , the sequence \((\alpha_n)\) is strictly increasing and the sequence \((\beta_n)\) is strictly decreasing; conclude using Set Theory, Chapter III, §4, Exercise 3.)
\item
  Let \(S(\Gamma, k)\) be the set of families \(x = (x_{\alpha})_{\alpha \in \Gamma}\) such that \(x_{\alpha} \in k\) for all \(\alpha \in \Gamma\) and the set of \(\alpha\) for which \(x_{\alpha} \neq 0\) is a well ordered subset of \(\Gamma\) . Show that \(S(\Gamma, k)\) is an additive subgroup of \(k^{\Gamma}\) ; moreover, if \(x = (x_{\alpha}), y = (y_{\alpha})\) are two elements of \(S(\Gamma, k)\) , the set \(C\) of \(\alpha + \beta\) for the ordered pairs \((\alpha, \beta)\) such that \(x_{\alpha} \neq 0\) and \(y_{\beta} \neq 0\) is well ordered by (a) and, for all \(\gamma \in C\) , the element \(z_{\gamma} = \sum_{\alpha + \beta = \gamma} x_{\alpha} y_{\beta}\) is defined in \(k\) ; if we write \(z_{\alpha} = 0\) for \(\alpha \notin C\) , the element \(z = (z_{\alpha}) \in k^{\Gamma}\) belongs to \(S(\Gamma, k)\) ; it is denoted by \(xy\) . Show that with this law of composition and addition on the product group \(k^{\Gamma}\) the set \(S(\Gamma, k)\) is a field; moreover, for all \(x = (x_{\alpha}) \neq 0\) of \(S(\Gamma, k)\) , let \(\lambda\) be the least of the \(\alpha \in \Gamma\) such that \(x_{\alpha} \neq 0\) ; if we write \(v(x) = \lambda\) (and \(v(0) = +\infty\) ), show that \(v\) is a valuation on \(S(\Gamma, k)\) with \(k\) as residue field and \(\Gamma\) as value group. \(v\) is called the canonical valuation on the field \(S(\Gamma, k)\) . The field \(k\) is canonically identified with a subfield of \(S(\Gamma, k)\) . For all \(\alpha \in \Gamma\) , let \(u_{\alpha}\) denote the element \((x_{\lambda})_{\lambda \in \Gamma}\) such that \(x_{\alpha} = 1, x_{\lambda} = 0\) for \(\lambda \neq \alpha\) ; for all non-zero \(z \in S(\Gamma, k)\) , there is then a unique ordered pair \((t, \alpha) \in k \times \Gamma\) such that \(v(z) = \alpha\) and \(v(z - tu_{\alpha}) > \alpha\) . \(tu_{\alpha}\) is called the dominant term of \(z\) .
\end{enumerate}

¶ 3. Let K be a field, w a valuation on K, Γ the value group of w and k the residue field of w. We form the corresponding field S(Γ, k) and preserve the notation of Exercise 2. For all \(t \in k\), let \(t^{0}\) be an element of the class t in the valuation ring of w; for all \(\alpha \in \Gamma\), let \(u_{\alpha}^{0}\) be an element of K such that \(w(u_{\alpha}^{0}) = \alpha\).

\begin{enumerate}
\def\labelenumi{(\alph{enumi})}
\tightlist
\item
  Let M be a subset of S(Γ, k) containing the elements \(tu_{\alpha}\) for every ordered pair \((t, \alpha) \in k \times \Gamma\) ; suppose that there exists a bijection \(x \mapsto x^{0}\) of M onto a subset \(M^{0}\) of K and that the following conditions are fulfilled: (1) \((tu_{\alpha})^{0} = t^{0}u_{\alpha}^{0}\) for every ordered pair \((t, \alpha) \in k \times \Gamma\) ; (2) if \(x = (x_{\alpha})\) belongs to M and \(C_{x}\) is the well-ordered subset of Γ consisting of the α such that \(x_{\alpha} \neq 0\) , then for every segment D of \(C_{x}\) the element \((x_{\lambda})_{\lambda \in D}\) belongs to M; (3) if x, y are two elements of M and \(tu_{\alpha}\) is the dominant term of x - y, then \(w(x^{0} - y^{0} - t^{0}u_{\alpha}^{0}) > w(x^{0} - y^{0})\) . If \(z'\) is an element of K not belonging to \(M^{0}\) , show that there exists an element \(z \notin M\) in S(Γ, k) such that the mapping which coincides with \(x \mapsto x^{0}\) on M and maps \(z'\) to z also satisfies the above conditions. (Note that, if \(w(z') = \alpha\) , there exists a \(t^{0}\) such that
\end{enumerate}

\[
w (z ^ {\prime} - t ^ {0} u _ {\alpha} ^ {0}) > \alpha ;
\]

write \(z_{\alpha} = tu_{\alpha}\) , then show by transfinite induction that the \(z_{\beta}\) of index \(\beta > \alpha\) may be determined so that \(z = (z_{\lambda})_{\lambda \in \Gamma}\) solves the problem.)

\begin{enumerate}
\def\labelenumi{(\alph{enumi})}
\setcounter{enumi}{1}
\tightlist
\item
  Deduce from (a) that \(\operatorname{Card}(\mathbf{K}) \leqslant \operatorname{Card} S(\Gamma, k)\) .
\end{enumerate}

\begin{enumerate}
\def\labelenumi{\arabic{enumi}.}
\setcounter{enumi}{3}
\tightlist
\item
  Let A be a local integral domain, K its field of fractions, m its maximal ideal and v a valuation on K whose ring B dominates A. Suppose that, either m is a finitely generated ideal, or v is a discrete valuation; then there exists in m an element x such that \(v(x) = \inf_{y \in \mathfrak{m}} v(y)\) . Let \(A'\) be the subring of K generated by
\end{enumerate}

A and the elements \(yx^{-1}\) , where y runs through m; let \(A_{1}\) be the local ring of \(A'\) relative to the prime ideal \(p \cap A'\) , where p is the ideal of the valuation v. Show that the ring \(A_{1}\) does not depend on the choice of the element x satisfying the above conditions ( \(A_{1}\) is called ``the prime quadratic transform of A in the direction of \(v'\) ); if A is Noetherian, so is \(A_{1}\) .

¶ 5. Let k be a field and K = k(X, Y) the field of rational functions in two indeterminates over k. We write \(x_{1} = X\) , \(y_{1} = Y\) and for \(n \geqslant 1\) define inductively the elements \(x_{n+1} = y_{n}\) and \(y_{n+1} = x_{n}y_{n}^{-1}\) of K.

\begin{enumerate}
\def\labelenumi{(\alph{enumi})}
\tightlist
\item
  Write \(\mathbf{A}_n = k[x_n, y_n]\) ; the sequence \((\mathbf{A}_n)\) is increasing and, if
\end{enumerate}

\[
\mathfrak {m} _ {n} = \mathrm{A} _ {n} x _ {n} + \mathrm{A} _ {n} y _ {n},
\]

\(m_{n}\) is a maximal ideal of \(A_{n}\) and \(m_{n} = A_{n} \cap m_{n+1}\) ; in the ring \(A = \bigcap_{n} A_{n}\) , \(m = \bigcap_{n} m_{n}\) is a maximal ideal.

\begin{enumerate}
\def\labelenumi{(\alph{enumi})}
\setcounter{enumi}{1}
\item
  Let \(v\) be the valuation on \(\mathbf{K}\) defined by \(v(\mathbf{X}) = \rho = (1 + \sqrt{5}) / 2\) , \(v(\mathbf{Y}) = 1\) ; let \(\mathbf{B}\) and \(\mathfrak{p}\) be the ring and ideal of \(v\) ; show that \(\mathfrak{m} = \mathfrak{p} \cap \mathbf{A}\) (note that \(\rho^2 - \rho - 1 = 0\) and use this fact to calculate \(v(y_n)\) ). Show that for a monomial \(\mathbf{X}^i\mathbf{Y}^j\) ( \(i, j\) positive or negative integers) to belong to \(\mathbf{A}\) , it is necessary and sufficient that \(i\left(\frac{1 + \sqrt{5}}{2}\right) + j > 0\) ; for every monomial \(z = \mathbf{X}^i\mathbf{Y}^j\) , therefore, either \(z \in \mathbf{A}\) or \(1/z \in \mathbf{A}\) .
\item
  Show that \(B = A_{m}\) (write an element \(t \in B\) in the form of a quotient of two polynomials in X and Y and show in each of them the monomials at which v takes the least value).
\item
  Show that, for all \(n\) , \((\mathbf{A}_{n+1})_{\mathfrak{m}_{n+1}}\) is the prime quadratic transform of \((\mathbf{A}_n)_{\mathfrak{m}_n}\) in the direction of \(v\) (Exercise 4).
\end{enumerate}

\begin{enumerate}
\def\labelenumi{\arabic{enumi}.}
\setcounter{enumi}{5}
\tightlist
\item
  \begin{enumerate}
  \def\labelenumii{(\alph{enumii})}
  \tightlist
  \item
    Let \(\mathbf{A}\) be a valuation ring. Extend to finitely generated A-modules the results of Algebra, Chapter VII, § 4, nos. 1 to 4 (theory of invariant factors).
  \end{enumerate}
\end{enumerate}

\begin{enumerate}
\def\labelenumi{(\alph{enumi})}
\setcounter{enumi}{1}
\item
  Let K be a field, v a valuation on K, A the ring of v and \(\Gamma\) the order group of v. Let \(U = (\alpha_{ij})\) be a square matrix of order n with elements in A and of determinant 0; show that there exists \(\lambda \in \Gamma\) such that, for every square matrix \(U' = (\alpha_{ij}')\) of order n with elements in A and satisfying the conditions \(v(\alpha_{ij}' - \alpha_{ij}) > \lambda\) for every ordered pair \((i,j)\) , the invariant factors of \(U'\) are the same as those of U.
\item
  Under the hypotheses of (b), generalize the following results of Algebra, Chapter IX, § 5, no. 1, Theorem and Exercise 1 and § 6, Exercise 10 (for the last, suppose that K is not of characteristic 2 and that \(v(2) = 0\) ).
\end{enumerate}

\begin{enumerate}
\def\labelenumi{\arabic{enumi}.}
\setcounter{enumi}{6}
\item
  Show that the integral domain B defined in Chapter II, § 3, Exercise 2 (b) is a local ring whose maximal ideal is principal but is not a valuation ring.
\item
  Show that, if a valuation ring A is strongly Laskerian (Chapter IV, § 2, Exercise 28), A is a field or a discrete valuation ring (use Exercise 29 of Chapter IV, § 2).
\end{enumerate}
